\documentclass[11pt]{article}

    \usepackage[breakable]{tcolorbox}
    \usepackage{parskip} % Stop auto-indenting (to mimic markdown behaviour)
    

    % Basic figure setup, for now with no caption control since it's done
    % automatically by Pandoc (which extracts ![](path) syntax from Markdown).
    \usepackage{graphicx}
    % Keep aspect ratio if custom image width or height is specified
    \setkeys{Gin}{keepaspectratio}
    % Maintain compatibility with old templates. Remove in nbconvert 6.0
    \let\Oldincludegraphics\includegraphics
    % Ensure that by default, figures have no caption (until we provide a
    % proper Figure object with a Caption API and a way to capture that
    % in the conversion process - todo).
    \usepackage{caption}
    \DeclareCaptionFormat{nocaption}{}
    \captionsetup{format=nocaption,aboveskip=0pt,belowskip=0pt}

    \usepackage{float}
    \floatplacement{figure}{H} % forces figures to be placed at the correct location
    \usepackage{xcolor} % Allow colors to be defined
    \usepackage{enumerate} % Needed for markdown enumerations to work
    \usepackage{geometry} % Used to adjust the document margins
    \usepackage{amsmath} % Equations
    \usepackage{amssymb} % Equations
    \usepackage{textcomp} % defines textquotesingle
    % Hack from http://tex.stackexchange.com/a/47451/13684:
    \AtBeginDocument{%
        \def\PYZsq{\textquotesingle}% Upright quotes in Pygmentized code
    }
    \usepackage{upquote} % Upright quotes for verbatim code
    \usepackage{eurosym} % defines \euro

    \usepackage{iftex}
    \ifPDFTeX
        \usepackage[T1]{fontenc}
        \IfFileExists{alphabeta.sty}{
              \usepackage{alphabeta}
          }{
              \usepackage[mathletters]{ucs}
              \usepackage[utf8x]{inputenc}
          }
    \else
        \usepackage{fontspec}
        \usepackage{unicode-math}
    \fi

    \usepackage{fancyvrb} % verbatim replacement that allows latex
    \usepackage{grffile} % extends the file name processing of package graphics
                         % to support a larger range
    \makeatletter % fix for old versions of grffile with XeLaTeX
    \@ifpackagelater{grffile}{2019/11/01}
    {
      % Do nothing on new versions
    }
    {
      \def\Gread@@xetex#1{%
        \IfFileExists{"\Gin@base".bb}%
        {\Gread@eps{\Gin@base.bb}}%
        {\Gread@@xetex@aux#1}%
      }
    }
    \makeatother
    \usepackage[Export]{adjustbox} % Used to constrain images to a maximum size
    \adjustboxset{max size={0.9\linewidth}{0.9\paperheight}}

    % The hyperref package gives us a pdf with properly built
    % internal navigation ('pdf bookmarks' for the table of contents,
    % internal cross-reference links, web links for URLs, etc.)
    \usepackage{hyperref}
    % The default LaTeX title has an obnoxious amount of whitespace. By default,
    % titling removes some of it. It also provides customization options.
    \usepackage{titling}
    \usepackage{longtable} % longtable support required by pandoc >1.10
    \usepackage{booktabs}  % table support for pandoc > 1.12.2
    \usepackage{array}     % table support for pandoc >= 2.11.3
    \usepackage{calc}      % table minipage width calculation for pandoc >= 2.11.1
    \usepackage[inline]{enumitem} % IRkernel/repr support (it uses the enumerate* environment)
    \usepackage[normalem]{ulem} % ulem is needed to support strikethroughs (\sout)
                                % normalem makes italics be italics, not underlines
    \usepackage{soul}      % strikethrough (\st) support for pandoc >= 3.0.0
    \usepackage{mathrsfs}
    

    
    % Colors for the hyperref package
    \definecolor{urlcolor}{rgb}{0,.145,.698}
    \definecolor{linkcolor}{rgb}{.71,0.21,0.01}
    \definecolor{citecolor}{rgb}{.12,.54,.11}

    % ANSI colors
    \definecolor{ansi-black}{HTML}{3E424D}
    \definecolor{ansi-black-intense}{HTML}{282C36}
    \definecolor{ansi-red}{HTML}{E75C58}
    \definecolor{ansi-red-intense}{HTML}{B22B31}
    \definecolor{ansi-green}{HTML}{00A250}
    \definecolor{ansi-green-intense}{HTML}{007427}
    \definecolor{ansi-yellow}{HTML}{DDB62B}
    \definecolor{ansi-yellow-intense}{HTML}{B27D12}
    \definecolor{ansi-blue}{HTML}{208FFB}
    \definecolor{ansi-blue-intense}{HTML}{0065CA}
    \definecolor{ansi-magenta}{HTML}{D160C4}
    \definecolor{ansi-magenta-intense}{HTML}{A03196}
    \definecolor{ansi-cyan}{HTML}{60C6C8}
    \definecolor{ansi-cyan-intense}{HTML}{258F8F}
    \definecolor{ansi-white}{HTML}{C5C1B4}
    \definecolor{ansi-white-intense}{HTML}{A1A6B2}
    \definecolor{ansi-default-inverse-fg}{HTML}{FFFFFF}
    \definecolor{ansi-default-inverse-bg}{HTML}{000000}

    % common color for the border for error outputs.
    \definecolor{outerrorbackground}{HTML}{FFDFDF}

    % commands and environments needed by pandoc snippets
    % extracted from the output of `pandoc -s`
    \providecommand{\tightlist}{%
      \setlength{\itemsep}{0pt}\setlength{\parskip}{0pt}}
    \DefineVerbatimEnvironment{Highlighting}{Verbatim}{commandchars=\\\{\}}
    % Add ',fontsize=\small' for more characters per line
    \newenvironment{Shaded}{}{}
    \newcommand{\KeywordTok}[1]{\textcolor[rgb]{0.00,0.44,0.13}{\textbf{{#1}}}}
    \newcommand{\DataTypeTok}[1]{\textcolor[rgb]{0.56,0.13,0.00}{{#1}}}
    \newcommand{\DecValTok}[1]{\textcolor[rgb]{0.25,0.63,0.44}{{#1}}}
    \newcommand{\BaseNTok}[1]{\textcolor[rgb]{0.25,0.63,0.44}{{#1}}}
    \newcommand{\FloatTok}[1]{\textcolor[rgb]{0.25,0.63,0.44}{{#1}}}
    \newcommand{\CharTok}[1]{\textcolor[rgb]{0.25,0.44,0.63}{{#1}}}
    \newcommand{\StringTok}[1]{\textcolor[rgb]{0.25,0.44,0.63}{{#1}}}
    \newcommand{\CommentTok}[1]{\textcolor[rgb]{0.38,0.63,0.69}{\textit{{#1}}}}
    \newcommand{\OtherTok}[1]{\textcolor[rgb]{0.00,0.44,0.13}{{#1}}}
    \newcommand{\AlertTok}[1]{\textcolor[rgb]{1.00,0.00,0.00}{\textbf{{#1}}}}
    \newcommand{\FunctionTok}[1]{\textcolor[rgb]{0.02,0.16,0.49}{{#1}}}
    \newcommand{\RegionMarkerTok}[1]{{#1}}
    \newcommand{\ErrorTok}[1]{\textcolor[rgb]{1.00,0.00,0.00}{\textbf{{#1}}}}
    \newcommand{\NormalTok}[1]{{#1}}

    % Additional commands for more recent versions of Pandoc
    \newcommand{\ConstantTok}[1]{\textcolor[rgb]{0.53,0.00,0.00}{{#1}}}
    \newcommand{\SpecialCharTok}[1]{\textcolor[rgb]{0.25,0.44,0.63}{{#1}}}
    \newcommand{\VerbatimStringTok}[1]{\textcolor[rgb]{0.25,0.44,0.63}{{#1}}}
    \newcommand{\SpecialStringTok}[1]{\textcolor[rgb]{0.73,0.40,0.53}{{#1}}}
    \newcommand{\ImportTok}[1]{{#1}}
    \newcommand{\DocumentationTok}[1]{\textcolor[rgb]{0.73,0.13,0.13}{\textit{{#1}}}}
    \newcommand{\AnnotationTok}[1]{\textcolor[rgb]{0.38,0.63,0.69}{\textbf{\textit{{#1}}}}}
    \newcommand{\CommentVarTok}[1]{\textcolor[rgb]{0.38,0.63,0.69}{\textbf{\textit{{#1}}}}}
    \newcommand{\VariableTok}[1]{\textcolor[rgb]{0.10,0.09,0.49}{{#1}}}
    \newcommand{\ControlFlowTok}[1]{\textcolor[rgb]{0.00,0.44,0.13}{\textbf{{#1}}}}
    \newcommand{\OperatorTok}[1]{\textcolor[rgb]{0.40,0.40,0.40}{{#1}}}
    \newcommand{\BuiltInTok}[1]{{#1}}
    \newcommand{\ExtensionTok}[1]{{#1}}
    \newcommand{\PreprocessorTok}[1]{\textcolor[rgb]{0.74,0.48,0.00}{{#1}}}
    \newcommand{\AttributeTok}[1]{\textcolor[rgb]{0.49,0.56,0.16}{{#1}}}
    \newcommand{\InformationTok}[1]{\textcolor[rgb]{0.38,0.63,0.69}{\textbf{\textit{{#1}}}}}
    \newcommand{\WarningTok}[1]{\textcolor[rgb]{0.38,0.63,0.69}{\textbf{\textit{{#1}}}}}
    \makeatletter
    \newsavebox\pandoc@box
    \newcommand*\pandocbounded[1]{%
      \sbox\pandoc@box{#1}%
      % scaling factors for width and height
      \Gscale@div\@tempa\textheight{\dimexpr\ht\pandoc@box+\dp\pandoc@box\relax}%
      \Gscale@div\@tempb\linewidth{\wd\pandoc@box}%
      % select the smaller of both
      \ifdim\@tempb\p@<\@tempa\p@
        \let\@tempa\@tempb
      \fi
      % scaling accordingly (\@tempa < 1)
      \ifdim\@tempa\p@<\p@
        \scalebox{\@tempa}{\usebox\pandoc@box}%
      % scaling not needed, use as it is
      \else
        \usebox{\pandoc@box}%
      \fi
    }
    \makeatother

    % Define a nice break command that doesn't care if a line doesn't already
    % exist.
    \def\br{\hspace*{\fill} \\* }
    % Math Jax compatibility definitions
    \def\gt{>}
    \def\lt{<}
    \let\Oldtex\TeX
    \let\Oldlatex\LaTeX
    \renewcommand{\TeX}{\textrm{\Oldtex}}
    \renewcommand{\LaTeX}{\textrm{\Oldlatex}}
    % Document parameters
    % Document title
    \title{Trabalho\_Trilha\_I\_AcolheMente\_PB}
    
    
    
    
    
    
    
% Pygments definitions
\makeatletter
\def\PY@reset{\let\PY@it=\relax \let\PY@bf=\relax%
    \let\PY@ul=\relax \let\PY@tc=\relax%
    \let\PY@bc=\relax \let\PY@ff=\relax}
\def\PY@tok#1{\csname PY@tok@#1\endcsname}
\def\PY@toks#1+{\ifx\relax#1\empty\else%
    \PY@tok{#1}\expandafter\PY@toks\fi}
\def\PY@do#1{\PY@bc{\PY@tc{\PY@ul{%
    \PY@it{\PY@bf{\PY@ff{#1}}}}}}}
\def\PY#1#2{\PY@reset\PY@toks#1+\relax+\PY@do{#2}}

\@namedef{PY@tok@w}{\def\PY@tc##1{\textcolor[rgb]{0.73,0.73,0.73}{##1}}}
\@namedef{PY@tok@c}{\let\PY@it=\textit\def\PY@tc##1{\textcolor[rgb]{0.24,0.48,0.48}{##1}}}
\@namedef{PY@tok@cp}{\def\PY@tc##1{\textcolor[rgb]{0.61,0.40,0.00}{##1}}}
\@namedef{PY@tok@k}{\let\PY@bf=\textbf\def\PY@tc##1{\textcolor[rgb]{0.00,0.50,0.00}{##1}}}
\@namedef{PY@tok@kp}{\def\PY@tc##1{\textcolor[rgb]{0.00,0.50,0.00}{##1}}}
\@namedef{PY@tok@kt}{\def\PY@tc##1{\textcolor[rgb]{0.69,0.00,0.25}{##1}}}
\@namedef{PY@tok@o}{\def\PY@tc##1{\textcolor[rgb]{0.40,0.40,0.40}{##1}}}
\@namedef{PY@tok@ow}{\let\PY@bf=\textbf\def\PY@tc##1{\textcolor[rgb]{0.67,0.13,1.00}{##1}}}
\@namedef{PY@tok@nb}{\def\PY@tc##1{\textcolor[rgb]{0.00,0.50,0.00}{##1}}}
\@namedef{PY@tok@nf}{\def\PY@tc##1{\textcolor[rgb]{0.00,0.00,1.00}{##1}}}
\@namedef{PY@tok@nc}{\let\PY@bf=\textbf\def\PY@tc##1{\textcolor[rgb]{0.00,0.00,1.00}{##1}}}
\@namedef{PY@tok@nn}{\let\PY@bf=\textbf\def\PY@tc##1{\textcolor[rgb]{0.00,0.00,1.00}{##1}}}
\@namedef{PY@tok@ne}{\let\PY@bf=\textbf\def\PY@tc##1{\textcolor[rgb]{0.80,0.25,0.22}{##1}}}
\@namedef{PY@tok@nv}{\def\PY@tc##1{\textcolor[rgb]{0.10,0.09,0.49}{##1}}}
\@namedef{PY@tok@no}{\def\PY@tc##1{\textcolor[rgb]{0.53,0.00,0.00}{##1}}}
\@namedef{PY@tok@nl}{\def\PY@tc##1{\textcolor[rgb]{0.46,0.46,0.00}{##1}}}
\@namedef{PY@tok@ni}{\let\PY@bf=\textbf\def\PY@tc##1{\textcolor[rgb]{0.44,0.44,0.44}{##1}}}
\@namedef{PY@tok@na}{\def\PY@tc##1{\textcolor[rgb]{0.41,0.47,0.13}{##1}}}
\@namedef{PY@tok@nt}{\let\PY@bf=\textbf\def\PY@tc##1{\textcolor[rgb]{0.00,0.50,0.00}{##1}}}
\@namedef{PY@tok@nd}{\def\PY@tc##1{\textcolor[rgb]{0.67,0.13,1.00}{##1}}}
\@namedef{PY@tok@s}{\def\PY@tc##1{\textcolor[rgb]{0.73,0.13,0.13}{##1}}}
\@namedef{PY@tok@sd}{\let\PY@it=\textit\def\PY@tc##1{\textcolor[rgb]{0.73,0.13,0.13}{##1}}}
\@namedef{PY@tok@si}{\let\PY@bf=\textbf\def\PY@tc##1{\textcolor[rgb]{0.64,0.35,0.47}{##1}}}
\@namedef{PY@tok@se}{\let\PY@bf=\textbf\def\PY@tc##1{\textcolor[rgb]{0.67,0.36,0.12}{##1}}}
\@namedef{PY@tok@sr}{\def\PY@tc##1{\textcolor[rgb]{0.64,0.35,0.47}{##1}}}
\@namedef{PY@tok@ss}{\def\PY@tc##1{\textcolor[rgb]{0.10,0.09,0.49}{##1}}}
\@namedef{PY@tok@sx}{\def\PY@tc##1{\textcolor[rgb]{0.00,0.50,0.00}{##1}}}
\@namedef{PY@tok@m}{\def\PY@tc##1{\textcolor[rgb]{0.40,0.40,0.40}{##1}}}
\@namedef{PY@tok@gh}{\let\PY@bf=\textbf\def\PY@tc##1{\textcolor[rgb]{0.00,0.00,0.50}{##1}}}
\@namedef{PY@tok@gu}{\let\PY@bf=\textbf\def\PY@tc##1{\textcolor[rgb]{0.50,0.00,0.50}{##1}}}
\@namedef{PY@tok@gd}{\def\PY@tc##1{\textcolor[rgb]{0.63,0.00,0.00}{##1}}}
\@namedef{PY@tok@gi}{\def\PY@tc##1{\textcolor[rgb]{0.00,0.52,0.00}{##1}}}
\@namedef{PY@tok@gr}{\def\PY@tc##1{\textcolor[rgb]{0.89,0.00,0.00}{##1}}}
\@namedef{PY@tok@ge}{\let\PY@it=\textit}
\@namedef{PY@tok@gs}{\let\PY@bf=\textbf}
\@namedef{PY@tok@ges}{\let\PY@bf=\textbf\let\PY@it=\textit}
\@namedef{PY@tok@gp}{\let\PY@bf=\textbf\def\PY@tc##1{\textcolor[rgb]{0.00,0.00,0.50}{##1}}}
\@namedef{PY@tok@go}{\def\PY@tc##1{\textcolor[rgb]{0.44,0.44,0.44}{##1}}}
\@namedef{PY@tok@gt}{\def\PY@tc##1{\textcolor[rgb]{0.00,0.27,0.87}{##1}}}
\@namedef{PY@tok@err}{\def\PY@bc##1{{\setlength{\fboxsep}{\string -\fboxrule}\fcolorbox[rgb]{1.00,0.00,0.00}{1,1,1}{\strut ##1}}}}
\@namedef{PY@tok@kc}{\let\PY@bf=\textbf\def\PY@tc##1{\textcolor[rgb]{0.00,0.50,0.00}{##1}}}
\@namedef{PY@tok@kd}{\let\PY@bf=\textbf\def\PY@tc##1{\textcolor[rgb]{0.00,0.50,0.00}{##1}}}
\@namedef{PY@tok@kn}{\let\PY@bf=\textbf\def\PY@tc##1{\textcolor[rgb]{0.00,0.50,0.00}{##1}}}
\@namedef{PY@tok@kr}{\let\PY@bf=\textbf\def\PY@tc##1{\textcolor[rgb]{0.00,0.50,0.00}{##1}}}
\@namedef{PY@tok@bp}{\def\PY@tc##1{\textcolor[rgb]{0.00,0.50,0.00}{##1}}}
\@namedef{PY@tok@fm}{\def\PY@tc##1{\textcolor[rgb]{0.00,0.00,1.00}{##1}}}
\@namedef{PY@tok@vc}{\def\PY@tc##1{\textcolor[rgb]{0.10,0.09,0.49}{##1}}}
\@namedef{PY@tok@vg}{\def\PY@tc##1{\textcolor[rgb]{0.10,0.09,0.49}{##1}}}
\@namedef{PY@tok@vi}{\def\PY@tc##1{\textcolor[rgb]{0.10,0.09,0.49}{##1}}}
\@namedef{PY@tok@vm}{\def\PY@tc##1{\textcolor[rgb]{0.10,0.09,0.49}{##1}}}
\@namedef{PY@tok@sa}{\def\PY@tc##1{\textcolor[rgb]{0.73,0.13,0.13}{##1}}}
\@namedef{PY@tok@sb}{\def\PY@tc##1{\textcolor[rgb]{0.73,0.13,0.13}{##1}}}
\@namedef{PY@tok@sc}{\def\PY@tc##1{\textcolor[rgb]{0.73,0.13,0.13}{##1}}}
\@namedef{PY@tok@dl}{\def\PY@tc##1{\textcolor[rgb]{0.73,0.13,0.13}{##1}}}
\@namedef{PY@tok@s2}{\def\PY@tc##1{\textcolor[rgb]{0.73,0.13,0.13}{##1}}}
\@namedef{PY@tok@sh}{\def\PY@tc##1{\textcolor[rgb]{0.73,0.13,0.13}{##1}}}
\@namedef{PY@tok@s1}{\def\PY@tc##1{\textcolor[rgb]{0.73,0.13,0.13}{##1}}}
\@namedef{PY@tok@mb}{\def\PY@tc##1{\textcolor[rgb]{0.40,0.40,0.40}{##1}}}
\@namedef{PY@tok@mf}{\def\PY@tc##1{\textcolor[rgb]{0.40,0.40,0.40}{##1}}}
\@namedef{PY@tok@mh}{\def\PY@tc##1{\textcolor[rgb]{0.40,0.40,0.40}{##1}}}
\@namedef{PY@tok@mi}{\def\PY@tc##1{\textcolor[rgb]{0.40,0.40,0.40}{##1}}}
\@namedef{PY@tok@il}{\def\PY@tc##1{\textcolor[rgb]{0.40,0.40,0.40}{##1}}}
\@namedef{PY@tok@mo}{\def\PY@tc##1{\textcolor[rgb]{0.40,0.40,0.40}{##1}}}
\@namedef{PY@tok@ch}{\let\PY@it=\textit\def\PY@tc##1{\textcolor[rgb]{0.24,0.48,0.48}{##1}}}
\@namedef{PY@tok@cm}{\let\PY@it=\textit\def\PY@tc##1{\textcolor[rgb]{0.24,0.48,0.48}{##1}}}
\@namedef{PY@tok@cpf}{\let\PY@it=\textit\def\PY@tc##1{\textcolor[rgb]{0.24,0.48,0.48}{##1}}}
\@namedef{PY@tok@c1}{\let\PY@it=\textit\def\PY@tc##1{\textcolor[rgb]{0.24,0.48,0.48}{##1}}}
\@namedef{PY@tok@cs}{\let\PY@it=\textit\def\PY@tc##1{\textcolor[rgb]{0.24,0.48,0.48}{##1}}}

\def\PYZbs{\char`\\}
\def\PYZus{\char`\_}
\def\PYZob{\char`\{}
\def\PYZcb{\char`\}}
\def\PYZca{\char`\^}
\def\PYZam{\char`\&}
\def\PYZlt{\char`\<}
\def\PYZgt{\char`\>}
\def\PYZsh{\char`\#}
\def\PYZpc{\char`\%}
\def\PYZdl{\char`\$}
\def\PYZhy{\char`\-}
\def\PYZsq{\char`\'}
\def\PYZdq{\char`\"}
\def\PYZti{\char`\~}
% for compatibility with earlier versions
\def\PYZat{@}
\def\PYZlb{[}
\def\PYZrb{]}
\makeatother


    % For linebreaks inside Verbatim environment from package fancyvrb.
    \makeatletter
        \newbox\Wrappedcontinuationbox
        \newbox\Wrappedvisiblespacebox
        \newcommand*\Wrappedvisiblespace {\textcolor{red}{\textvisiblespace}}
        \newcommand*\Wrappedcontinuationsymbol {\textcolor{red}{\llap{\tiny$\m@th\hookrightarrow$}}}
        \newcommand*\Wrappedcontinuationindent {3ex }
        \newcommand*\Wrappedafterbreak {\kern\Wrappedcontinuationindent\copy\Wrappedcontinuationbox}
        % Take advantage of the already applied Pygments mark-up to insert
        % potential linebreaks for TeX processing.
        %        {, <, #, %, $, ' and ": go to next line.
        %        _, }, ^, &, >, - and ~: stay at end of broken line.
        % Use of \textquotesingle for straight quote.
        \newcommand*\Wrappedbreaksatspecials {%
            \def\PYGZus{\discretionary{\char`\_}{\Wrappedafterbreak}{\char`\_}}%
            \def\PYGZob{\discretionary{}{\Wrappedafterbreak\char`\{}{\char`\{}}%
            \def\PYGZcb{\discretionary{\char`\}}{\Wrappedafterbreak}{\char`\}}}%
            \def\PYGZca{\discretionary{\char`\^}{\Wrappedafterbreak}{\char`\^}}%
            \def\PYGZam{\discretionary{\char`\&}{\Wrappedafterbreak}{\char`\&}}%
            \def\PYGZlt{\discretionary{}{\Wrappedafterbreak\char`\<}{\char`\<}}%
            \def\PYGZgt{\discretionary{\char`\>}{\Wrappedafterbreak}{\char`\>}}%
            \def\PYGZsh{\discretionary{}{\Wrappedafterbreak\char`\#}{\char`\#}}%
            \def\PYGZpc{\discretionary{}{\Wrappedafterbreak\char`\%}{\char`\%}}%
            \def\PYGZdl{\discretionary{}{\Wrappedafterbreak\char`\$}{\char`\$}}%
            \def\PYGZhy{\discretionary{\char`\-}{\Wrappedafterbreak}{\char`\-}}%
            \def\PYGZsq{\discretionary{}{\Wrappedafterbreak\textquotesingle}{\textquotesingle}}%
            \def\PYGZdq{\discretionary{}{\Wrappedafterbreak\char`\"}{\char`\"}}%
            \def\PYGZti{\discretionary{\char`\~}{\Wrappedafterbreak}{\char`\~}}%
        }
        % Some characters . , ; ? ! / are not pygmentized.
        % This macro makes them "active" and they will insert potential linebreaks
        \newcommand*\Wrappedbreaksatpunct {%
            \lccode`\~`\.\lowercase{\def~}{\discretionary{\hbox{\char`\.}}{\Wrappedafterbreak}{\hbox{\char`\.}}}%
            \lccode`\~`\,\lowercase{\def~}{\discretionary{\hbox{\char`\,}}{\Wrappedafterbreak}{\hbox{\char`\,}}}%
            \lccode`\~`\;\lowercase{\def~}{\discretionary{\hbox{\char`\;}}{\Wrappedafterbreak}{\hbox{\char`\;}}}%
            \lccode`\~`\:\lowercase{\def~}{\discretionary{\hbox{\char`\:}}{\Wrappedafterbreak}{\hbox{\char`\:}}}%
            \lccode`\~`\?\lowercase{\def~}{\discretionary{\hbox{\char`\?}}{\Wrappedafterbreak}{\hbox{\char`\?}}}%
            \lccode`\~`\!\lowercase{\def~}{\discretionary{\hbox{\char`\!}}{\Wrappedafterbreak}{\hbox{\char`\!}}}%
            \lccode`\~`\/\lowercase{\def~}{\discretionary{\hbox{\char`\/}}{\Wrappedafterbreak}{\hbox{\char`\/}}}%
            \catcode`\.\active
            \catcode`\,\active
            \catcode`\;\active
            \catcode`\:\active
            \catcode`\?\active
            \catcode`\!\active
            \catcode`\/\active
            \lccode`\~`\~
        }
    \makeatother

    \let\OriginalVerbatim=\Verbatim
    \makeatletter
    \renewcommand{\Verbatim}[1][1]{%
        %\parskip\z@skip
        \sbox\Wrappedcontinuationbox {\Wrappedcontinuationsymbol}%
        \sbox\Wrappedvisiblespacebox {\FV@SetupFont\Wrappedvisiblespace}%
        \def\FancyVerbFormatLine ##1{\hsize\linewidth
            \vtop{\raggedright\hyphenpenalty\z@\exhyphenpenalty\z@
                \doublehyphendemerits\z@\finalhyphendemerits\z@
                \strut ##1\strut}%
        }%
        % If the linebreak is at a space, the latter will be displayed as visible
        % space at end of first line, and a continuation symbol starts next line.
        % Stretch/shrink are however usually zero for typewriter font.
        \def\FV@Space {%
            \nobreak\hskip\z@ plus\fontdimen3\font minus\fontdimen4\font
            \discretionary{\copy\Wrappedvisiblespacebox}{\Wrappedafterbreak}
            {\kern\fontdimen2\font}%
        }%

        % Allow breaks at special characters using \PYG... macros.
        \Wrappedbreaksatspecials
        % Breaks at punctuation characters . , ; ? ! and / need catcode=\active
        \OriginalVerbatim[#1,codes*=\Wrappedbreaksatpunct]%
    }
    \makeatother

    % Exact colors from NB
    \definecolor{incolor}{HTML}{303F9F}
    \definecolor{outcolor}{HTML}{D84315}
    \definecolor{cellborder}{HTML}{CFCFCF}
    \definecolor{cellbackground}{HTML}{F7F7F7}

    % prompt
    \makeatletter
    \newcommand{\boxspacing}{\kern\kvtcb@left@rule\kern\kvtcb@boxsep}
    \makeatother
    \newcommand{\prompt}[4]{
        {\ttfamily\llap{{\color{#2}[#3]:\hspace{3pt}#4}}\vspace{-\baselineskip}}
    }
    

    
    % Prevent overflowing lines due to hard-to-break entities
    \sloppy
    % Setup hyperref package
    \hypersetup{
      breaklinks=true,  % so long urls are correctly broken across lines
      colorlinks=true,
      urlcolor=urlcolor,
      linkcolor=linkcolor,
      citecolor=citecolor,
      }
    % Slightly bigger margins than the latex defaults
    
    \geometry{verbose,tmargin=1in,bmargin=1in,lmargin=1in,rmargin=1in}
    
    

\begin{document}
    
    \maketitle
    
    

    
    \begin{tcolorbox}[breakable, size=fbox, boxrule=1pt, pad at break*=1mm,colback=cellbackground, colframe=cellborder]
\prompt{In}{incolor}{ }{\boxspacing}
\begin{Verbatim}[commandchars=\\\{\}]
\PY{k+kn}{import}\PY{+w}{ }\PY{n+nn}{sys}\PY{o}{,}\PY{+w}{ }\PY{n+nn}{os}\PY{o}{,}\PY{+w}{ }\PY{n+nn}{shutil}\PY{o}{,}\PY{+w}{ }\PY{n+nn}{urllib}\PY{n+nn}{.}\PY{n+nn}{request}\PY{o}{,}\PY{+w}{ }\PY{n+nn}{warnings}
\PY{n}{warnings}\PY{o}{.}\PY{n}{filterwarnings}\PY{p}{(}\PY{l+s+s2}{\PYZdq{}}\PY{l+s+s2}{ignore}\PY{l+s+s2}{\PYZdq{}}\PY{p}{)}

\PY{c+c1}{\PYZsh{} Logo do IF Goiano}
\PY{n}{LOGO\PYZus{}URL} \PY{o}{=} \PY{p}{(}
    \PY{l+s+s2}{\PYZdq{}}\PY{l+s+s2}{https://www.ifgoiano.edu.br/home/images/}\PY{l+s+s2}{\PYZdq{}}
    \PY{l+s+s2}{\PYZdq{}}\PY{l+s+s2}{REITORIA/Imagens/2017/Jan\PYZus{}2017/logo\PYZus{}urutai.png}\PY{l+s+s2}{\PYZdq{}}
\PY{p}{)}
\PY{n}{LOGO\PYZus{}LOCAL} \PY{o}{=} \PY{l+s+s2}{\PYZdq{}}\PY{l+s+s2}{logo\PYZus{}urutai.png}\PY{l+s+s2}{\PYZdq{}}
\PY{k}{if} \PY{o+ow}{not} \PY{n}{os}\PY{o}{.}\PY{n}{path}\PY{o}{.}\PY{n}{exists}\PY{p}{(}\PY{n}{LOGO\PYZus{}LOCAL}\PY{p}{)}\PY{p}{:}
    \PY{k}{try}\PY{p}{:}
        \PY{n}{urllib}\PY{o}{.}\PY{n}{request}\PY{o}{.}\PY{n}{urlretrieve}\PY{p}{(}\PY{n}{LOGO\PYZus{}URL}\PY{p}{,} \PY{n}{LOGO\PYZus{}LOCAL}\PY{p}{)}
    \PY{k}{except} \PY{n+ne}{Exception}\PY{p}{:}
        \PY{k}{pass}

\PY{c+c1}{\PYZsh{} garante que o motor lógico é importável}
\PY{k}{if} \PY{l+s+s2}{\PYZdq{}}\PY{l+s+s2}{src}\PY{l+s+s2}{\PYZdq{}} \PY{o+ow}{not} \PY{o+ow}{in} \PY{n}{sys}\PY{o}{.}\PY{n}{path}\PY{p}{:}
    \PY{n}{sys}\PY{o}{.}\PY{n}{path}\PY{o}{.}\PY{n}{insert}\PY{p}{(}\PY{l+m+mi}{0}\PY{p}{,} \PY{n}{os}\PY{o}{.}\PY{n}{path}\PY{o}{.}\PY{n}{abspath}\PY{p}{(}
        \PY{n}{os}\PY{o}{.}\PY{n}{path}\PY{o}{.}\PY{n}{join}\PY{p}{(}\PY{n}{os}\PY{o}{.}\PY{n}{path}\PY{o}{.}\PY{n}{dirname}\PY{p}{(}\PY{l+s+s2}{\PYZdq{}}\PY{l+s+s2}{\PYZus{}\PYZus{}file\PYZus{}\PYZus{}}\PY{l+s+s2}{\PYZdq{}}\PY{p}{)}\PY{p}{,} \PY{l+s+s2}{\PYZdq{}}\PY{l+s+s2}{..}\PY{l+s+s2}{\PYZdq{}}\PY{p}{)}
    \PY{p}{)}\PY{p}{)}
\end{Verbatim}
\end{tcolorbox}

    \begin{titlepage}
    \centering

    \includegraphics[width=0.15\textwidth]{logo_urutai.png} \\[0.3cm]
    \vspace{0.8cm}

    {\small \textbf{INSTITUTO FEDERAL GOIANO — IF Goiano} \\}
    {\footnotesize PÓS-GRADUAÇÃO \textit{LATO SENSU} EM INTELIGÊNCIA ARTIFICIAL APLICADA \\}
    {\footnotesize TRILHA I — FUNDAMENTOS SIMBÓLICOS E BUSCA \\}
    \vspace{2.5cm}

    {\Large \MakeUppercase{Leonardo Maximino Bernardo} \par}
    \vspace{4cm}

    {\Large \textbf{AcolheMente Escolar PB: Motor de Inferência Lógico para Apoio à Gestão Escolar na Priorização Responsável de Acolhimento em Saúde Mental} \par}

    \vfill

    {\large Maio de 2026 \par}
\end{titlepage}
\newpage

    \section{1. Apresentação do Aluno e
Contexto}\label{apresentauxe7uxe3o-do-aluno-e-contexto}

Leonardo Maximino Bernardo é professor de Física da rede pública
estadual da Paraíba desde 2020, com mais de 15 anos de experiência em
docência nos níveis médio e superior. Doutor em Ciência e Engenharia de
Materiais pela Universidade Federal da Paraíba (UFPB), sua trajetória
acadêmica combina modelagem computacional avançada --- com uso de C++ e
OpenFOAM para simulação numérica de solidificação de ligas metálicas ---
e uma crescente atuação em ciência de dados, inteligência artificial e
desenvolvimento de software.

Ao longo de sua formação, Leonardo atuou como pesquisador na UFPB,
desenvolvendo rotinas de cálculo e algoritmos para análise de grandes
volumes de dados de simulação, extraindo padrões quantitativos para
publicações acadêmicas. Essa experiência consolidou habilidades de
pensamento analítico, resolução de problemas complexos e comunicação
científica que hoje informam diretamente sua prática docente e sua
abordagem em projetos de IA.

Atualmente, além da docência em Física, Leonardo cursa Pós-Graduação
\emph{Lato Sensu} em Inteligência Artificial Aplicada pelo Instituto
Federal Goiano (IF Goiano) e graduação em Ciência de Dados pelo Unipê.
Sua stack tecnológica inclui Python, JavaScript, SQL, frameworks como
Django e React, e ferramentas de análise como Pandas, Scikit-learn e
NetworkX. Possui certificações em desenvolvimento back-end com Node.js,
Harvard CS50 e aceleração global de desenvolvimento de software.

O presente trabalho nasce da intersecção entre sua vivência como
professor de escola pública estadual na Paraíba --- onde testemunha
diariamente demandas crescentes de acolhimento em saúde mental de
adolescentes --- e seu domínio técnico em modelagem computacional e
lógica formal. A escolha do domínio de saúde mental escolar não é
acidental: reflete o compromisso de aplicar inteligência artificial de
forma ética, explicável e governada em um contexto de vulnerabilidade
social, onde erros tecnológicos podem ter consequências humanas graves.

    \section{2. Introdução e Definição do
Problema}\label{introduuxe7uxe3o-e-definiuxe7uxe3o-do-problema}

A saúde mental de adolescentes no Brasil atravessa um momento crítico.
Dados da Pesquisa Nacional de Saúde do Escolar --- PeNSE 2024 (IBGE,
2024) revelam que parcela significativa dos estudantes de 13 a 17 anos
reporta sentimentos persistentes de tristeza, solidão e desesperança,
com índices particularmente preocupantes na região Nordeste e,
especificamente, no estado da Paraíba. O relatório da Organização
Mundial da Saúde (WHO, 2022) confirma essa tendência global: transtornos
mentais respondem por uma proporção crescente da carga de doença entre
jovens, e a maioria dos casos permanece sem qualquer forma de
acolhimento adequado.

No contexto da escola pública estadual paraibana, a demanda por
acolhimento em saúde mental frequentemente supera a capacidade de
resposta da equipe pedagógica e psicossocial. Professores, coordenadores
e orientadores educacionais enfrentam o desafio de identificar, dentre
centenas de estudantes, aqueles que necessitam de atenção prioritária
--- sem dispor de ferramentas estruturadas para essa tarefa. O resultado
é uma dependência excessiva de percepções individuais, sujeitas a viés
cognitivo, fadiga decisória e inconsistência entre turnos e unidades
escolares. Essa lacuna não é trivial: a ausência de priorização
sistemática pode significar que um adolescente em situação de risco
grave não receba acolhimento em tempo hábil.

Ao mesmo tempo, a adoção acrítica de tecnologias de inteligência
artificial nesse domínio apresenta riscos igualmente sérios. Sistemas
baseados em aprendizado de máquina (\emph{machine learning}) ou IA
generativa, embora poderosos em outros contextos, introduzem opacidade
decisória incompatível com o rigor ético exigido quando se trata de
menores de idade em situação de vulnerabilidade (Jobin \emph{et al.},
2019). A Lei Geral de Proteção de Dados --- LGPD (Brasil, 2018) e o
Estatuto da Criança e do Adolescente --- ECA (Brasil, 1990) impõem
salvaguardas rigorosas para o tratamento de dados sensíveis de menores,
incluindo o direito à explicação de decisões automatizadas. Um modelo de
``caixa preta'' que classifique ou ranqueie estudantes por risco
psicológico seria, no mínimo, juridicamente temerário e, no limite,
eticamente inaceitável.

Diante desse cenário, o presente trabalho propõe o \textbf{AcolheMente
Escolar PB}: um Sistema Baseado em Conhecimento (SBC) que utiliza
\textbf{Lógica Proposicional com Inferência por Resolução} para apoiar
--- e nunca substituir --- a gestão escolar na priorização responsável
de acolhimento em saúde mental. A escolha por lógica simbólica
determinística não é uma limitação técnica, mas uma decisão de
\emph{design} ético: garante explicabilidade total, rastreabilidade de
cada conclusão e ausência de qualquer linguagem diagnóstica clínica
(Russell; Norvig, 2022).

\subsubsection{Definição Operacional do
Problema}\label{definiuxe7uxe3o-operacional-do-problema}

\textbf{Questão central:} Como apoiar a gestão escolar de escolas
públicas estaduais da Paraíba na priorização sistemática de acolhimento
em saúde mental de adolescentes, utilizando inteligência artificial
explicável, sem produzir diagnósticos clínicos, sem violar a LGPD ou o
ECA, e sem substituir o julgamento humano profissional?

\textbf{Escopo:} Construção de um motor de inferência lógico
determinístico com 7 variáveis proposicionais e 6 regras em CNF,
alimentado por indicadores agregados da PeNSE 2024 e validado com
cenários sintéticos.

\textbf{Não escopo:} O sistema \textbf{NÃO diagnostica}, \textbf{NÃO
classifica} alunos individualmente, \textbf{NÃO} utiliza dados
identificáveis, \textbf{NÃO} substitui avaliação profissional de
psicólogos ou assistentes sociais, e \textbf{NÃO} se conecta a APIs de
IA generativa.

\textbf{Objetivo geral:} Desenvolver e validar um motor de inferência
por resolução, baseado em lógica proposicional e fundamentado em dados
oficiais da PeNSE 2024, para apoiar a priorização de acolhimento em
saúde mental escolar no estado da Paraíba.

\textbf{Objetivos específicos:}

\begin{enumerate}
\def\labelenumi{\arabic{enumi}.}
\tightlist
\item
  Modelar o domínio de saúde mental escolar em 7 variáveis
  proposicionais (5 nucleares + 2 contextuais), ancoradas em variáveis
  da PeNSE 2024.
\item
  Formalizar 6 regras de inferência em Forma Normal Conjuntiva (CNF) e
  implementar motor de resolução determinístico.
\item
  Validar o motor com cenários sintéticos representativos e comparar com
  \emph{baseline} manual.
\item
  Garantir conformidade integral com LGPD, ECA e Programa Saúde na
  Escola (PSE) (Brasil, 2007).
\item
  Documentar explicabilidade, limitações e perspectivas de evolução com
  rigor acadêmico.
\end{enumerate}

    \section{3. Metodologia e Escolha da Técnica de
IA}\label{metodologia-e-escolha-da-tuxe9cnica-de-ia}

\subsection{3.1 Sistemas Baseados em
Conhecimento}\label{sistemas-baseados-em-conhecimento}

O AcolheMente Escolar PB é, em sua essência, um \textbf{Sistema Baseado
em Conhecimento} (SBC) --- uma classe de sistemas de inteligência
artificial que codifica explicitamente o conhecimento de especialistas
humanos em regras formais e utiliza mecanismos de inferência para
derivar conclusões a partir de fatos observados (Brachman; Levesque,
2004). Diferentemente de abordagens estatísticas ou conexionistas, um
SBC não ``aprende'' a partir de dados brutos: ele opera sobre uma base
de conhecimento construída deliberadamente por engenheiros do
conhecimento em colaboração com especialistas do domínio.

Essa característica é particularmente valiosa no contexto de saúde
mental escolar, onde: (a) o volume de dados individuais disponível é
insuficiente para treinamento de modelos de \emph{machine learning}; (b)
a transparência de cada decisão é requisito ético e legal; e (c) o custo
de um erro --- falso positivo que rotula indevidamente um estudante, ou
falso negativo que ignora um caso grave --- é inaceitavelmente alto.

\subsection{3.2 Lógica Proposicional e Inferência por
Resolução}\label{luxf3gica-proposicional-e-inferuxeancia-por-resoluuxe7uxe3o}

A técnica específica adotada é a \textbf{Lógica Proposicional com
Inferência por Resolução}. Trata-se de um formalismo da lógica clássica
no qual o conhecimento é representado por proposições atômicas
(variáveis que assumem valores verdadeiro ou falso) e conectivos lógicos
(conjunção, disjunção, negação, implicação). A inferência é realizada
por meio do \textbf{método de resolução}, um procedimento correto e
refutacionalmente completo para a lógica proposicional (Genesereth;
Nilsson, 1987).

O método de resolução opera sobre cláusulas em \textbf{Forma Normal
Conjuntiva (CNF)} --- uma conjunção de disjunções de literais. Para
verificar se uma conclusão \(A\) pode ser derivada de uma base de
conhecimento \(KB\), o procedimento: (i) converte todas as regras para
CNF; (ii) adiciona a negação da conclusão (\(\neg A\)) ao conjunto de
cláusulas; (iii) aplica repetidamente a regra de resolução, buscando
derivar a cláusula vazia (\(\square\)), o que constitui prova por
contradição de que \(KB \models A\).

\subsection{3.3 Justificativa da Escolha e Alternativas
Descartadas}\label{justificativa-da-escolha-e-alternativas-descartadas}

A Trilha I do curso de Pós-Graduação em IA Aplicada do IF Goiano tem
como foco técnicas simbólicas e de busca. A escolha por lógica
proposicional com resolução é, portanto, duplamente justificada: atende
ao escopo pedagógico da trilha e oferece as propriedades de
explicabilidade e determinismo exigidas pelo domínio.

Foram consideradas e descartadas as seguintes alternativas:

\begin{itemize}
\tightlist
\item
  \textbf{Lógica de Primeira Ordem (FOL):} Embora mais expressiva,
  introduziria complexidade desnecessária para um domínio com 7
  variáveis binárias. A decidibilidade garantida da lógica proposicional
  é uma vantagem operacional.
\item
  \textbf{Redes Bayesianas:} Exigiriam distribuições de probabilidade
  condicional que não podem ser estimadas com rigor a partir de dados
  agregados da PeNSE. Além disso, a saída probabilística (``70\% de
  chance de necessitar acolhimento'') seria inadequada para um sistema
  que deve emitir apenas sinais binários de priorização.
\item
  \textbf{Árvores de Decisão / Random Forest:} Embora interpretáveis,
  dependem de treinamento supervisionado com dados rotulados ---
  inexistentes neste domínio. A criação de rótulos artificiais
  introduziria viés de confirmação.
\item
  \textbf{Redes Neurais / Deep Learning:} Incompatíveis com os
  requisitos de explicabilidade total e transparência jurídica. Como
  discutido na Seção 15, essas técnicas são admitidas como possibilidade
  teórica futura, condicionada a requisitos éticos e institucionais
  rigorosos.
\item
  \textbf{IA Generativa (LLMs):} Descartada por completo. Modelos de
  linguagem não oferecem garantias de consistência lógica, não são
  determinísticos, e sua utilização para classificação de risco em saúde
  mental de menores seria juridicamente indefensável sob a LGPD e o ECA.
\end{itemize}

Russell e Norvig (2022) argumentam que a escolha da representação do
conhecimento deve ser guiada pelas propriedades do domínio, e não pela
sofisticação da técnica. No caso do AcolheMente Escolar PB, a lógica
proposicional é a ferramenta \emph{correta}, não uma ferramenta
\emph{limitada}.

    \section{4. Dados e Governança}\label{dados-e-governanuxe7a}

\subsection{4.1 Fontes de Dados e Arquitetura de
Proveniência}\label{fontes-de-dados-e-arquitetura-de-proveniuxeancia}

O AcolheMente Escolar PB opera com uma arquitetura rigorosa de
proveniência de dados, dividida em três camadas (\emph{tiers})
mutuamente exclusivas. Essa separação não é apenas uma decisão técnica:
é um requisito de governança que impede contaminação cruzada entre
fontes de naturezas distintas.

\textbf{TIER\_A --- PeNSE 2024 (IBGE):} Base de dados oficial, de acesso
público, contendo indicadores agregados de saúde do escolar coletados
pelo Instituto Brasileiro de Geografia e Estatística (IBGE, 2024). A
PeNSE é uma pesquisa amostral com representatividade nacional, regional
e por unidade da federação. Os microdados utilizados referem-se
exclusivamente ao Tema 12 (Saúde Mental) e ao Tema 01 (Informações
Gerais), com recorte para o estado da Paraíba. É fundamental ressaltar
que a PeNSE não constitui prontuário individual: trata-se de dados
agregados, anonimizados na origem, sem possibilidade de identificação de
respondentes.

\textbf{TIER\_B --- Escola Sintética:} Conjunto de cenários gerados
computacionalmente para demonstração operacional e validação do motor
lógico. Cada cenário é um vetor binário de 7 posições (E, B, V, S, C, I,
A), representando combinações hipotéticas de indicadores. Nenhum dado
real de aluno ou escola é utilizado. A escola sintética permite
verificar a correção lógica do motor sem qualquer risco de privacidade.

\textbf{TIER\_C --- Dataset Externo (Benchmark):} Base de dados externa
utilizada exclusivamente para benchmark metodológico --- isto é, para
validar que o \emph{pipeline} de processamento funciona corretamente com
dados tabulares de outra origem. O dataset externo \textbf{não gera
conclusões sobre o Brasil, o Nordeste ou a Paraíba}. Qualquer tentativa
de inferência regional a partir do TIER\_C é logicamente inválida e
explicitamente proibida na arquitetura do sistema.

\textbf{Merge entre tiers é estritamente proibido.} A separação garante
que conclusões derivadas de dados oficiais (TIER\_A) nunca sejam
contaminadas por dados sintéticos (TIER\_B) ou externos (TIER\_C), e
vice-versa.

\subsection{4.2 Conformidade LGPD, ECA e
PSE}\label{conformidade-lgpd-eca-e-pse}

A Lei Geral de Proteção de Dados --- LGPD (Brasil, 2018) --- classifica
dados de saúde de menores de idade como dados pessoais sensíveis,
impondo requisitos de tratamento particularmente rigorosos: base legal
específica, finalidade determinada, minimização de dados e direito do
titular à explicação de decisões automatizadas. O AcolheMente Escolar PB
atende a esses requisitos \emph{by design}: (a) não processa dados
pessoais identificáveis; (b) opera sobre indicadores agregados
anonimizados; (c) toda conclusão é rastreável a regras explícitas; (d)
não produz decisões automatizadas sobre indivíduos.

O Estatuto da Criança e do Adolescente --- ECA (Brasil, 1990) --- veda
qualquer forma de rotulagem, classificação ou estigmatização de crianças
e adolescentes. O motor lógico do AcolheMente incorpora
\emph{guardrails} que impedem: uso de linguagem diagnóstica clínica
(depressão, TDAH, autismo, TOC, transtorno); ranking ou escore
comparativo entre estudantes; e classificação de alunos individuais
reais.

O Programa Saúde na Escola --- PSE (Brasil, 2007) --- estabelece a
integração entre escola e rede de saúde como política pública. O
AcolheMente não substitui essa integração: ele a apoia, oferecendo à
equipe escolar um instrumento de priorização que indica \emph{quando}
acionar os fluxos de encaminhamento previstos no PSE --- para a Unidade
Básica de Saúde (UBS), o Centro de Atenção Psicossocial (CAPS) ou o
Conselho Tutelar.

\subsection{4.3 Human-in-the-Loop}\label{human-in-the-loop}

O princípio de \emph{human-in-the-loop} (HITL) é inegociável no
AcolheMente. A variável de saída \textbf{A} (Acolhimento Priorizado) é
um \textbf{sinal para o profissional humano}, não uma decisão final. O
fluxo operacional previsto é: (1) o motor sinaliza priorização; (2) a
equipe pedagógica/psicossocial avalia o caso; (3) a decisão de ação é
sempre humana. Nenhum estudante será encaminhado, convocado ou
classificado com base exclusiva na saída do motor.

    \section{5. Representação do Conhecimento: Variáveis
Proposicionais}\label{representauxe7uxe3o-do-conhecimento-variuxe1veis-proposicionais}

\subsection{5.1 Engenharia do
Conhecimento}\label{engenharia-do-conhecimento}

A modelagem do domínio seguiu princípios clássicos de engenharia do
conhecimento (Brachman; Levesque, 2004): identificação de conceitos
relevantes, elicitação de regras com especialistas, formalização em
lógica proposicional e validação iterativa. O resultado é um conjunto de
7 variáveis proposicionais --- 5 nucleares e 2 contextuais --- que
representam o espaço de estados relevante para a priorização de
acolhimento.

A decisão de utilizar exatamente 7 variáveis --- e não 5 ou 10 --- foi
deliberada. Com 5 variáveis, perderíamos a capacidade de modelar fatores
contextuais (comportamento e resposta institucional) que modulam a
urgência do acolhimento. Com mais de 7, o espaço combinatório de \(2^n\)
estados cresceria sem ganho proporcional de discriminação, violando o
princípio de parcimônia que governa sistemas simbólicos interpretáveis.

\subsection{5.2 Variáveis Nucleares (5)}\label{variuxe1veis-nucleares-5}

\begin{itemize}
\tightlist
\item
  \textbf{E} --- Sofrimento emocional recorrente (choro, isolamento
  agudo). Fontes PeNSE: B12004, B12005, B12007.
\item
  \textbf{B} --- Baixo apoio familiar/social percebido. Fontes PeNSE:
  B12003, B07004.
\item
  \textbf{V} --- Indicador crítico de desvalor/desesperança em relação à
  vida. Fonte PeNSE: B12008.
\item
  \textbf{S} --- Sinal autorreferido de autoagressão. Fonte PeNSE:
  B12009.
\item
  \textbf{A} --- Acolhimento humano prioritário (\textbf{saída do
  motor}). Derivada por inferência.
\end{itemize}

\subsection{5.3 Variáveis Contextuais de Modulação
(2)}\label{variuxe1veis-contextuais-de-modulauxe7uxe3o-2}

\begin{itemize}
\tightlist
\item
  \textbf{C} --- Contexto comportamental agravante (mudança brusca de
  comportamento). Fontes PeNSE: B03010C, B03006B.
\item
  \textbf{I} --- Insuficiência institucional de resposta. Fontes PeNSE:
  E01P60, E01P117.
\end{itemize}

As variáveis contextuais \textbf{nunca} inferem \textbf{A} isoladamente.
\textbf{C} só participa de regras em conjunção com \textbf{E} e
\textbf{B} (R5); \textbf{I} só participa em conjunção com \textbf{V}
(R6). Essa restrição é um \emph{guardrail} lógico: impede que fatores
contextuais, por si sós, gerem falsos positivos de priorização.

\subsection{5.4 Nota sobre Não
Diagnóstico}\label{nota-sobre-nuxe3o-diagnuxf3stico}

O sistema \textbf{NÃO diagnostica} condições clínicas. As variáveis
representam indicadores comportamentais observáveis e auto-reportados,
não categorias nosológicas. A variável \textbf{A} sinaliza necessidade
de acolhimento pedagógico/psicossocial --- uma ação escolar, não uma
intervenção clínica. Essa distinção é fundamental e perpassa todo o
\emph{design} do motor.

    \section{6. Base de Conhecimento e Regras
Lógicas}\label{base-de-conhecimento-e-regras-luxf3gicas}

\subsection{6.1 Apresentação das Regras
R1--R6}\label{apresentauxe7uxe3o-das-regras-r1r6}

A base de conhecimento do AcolheMente Escolar PB é composta por 6 regras
de inferência, cada uma representando um cenário em que a combinação de
indicadores justifica a priorização de acolhimento. As regras foram
definidas com base em princípios pedagógicos e psicossociais, não em
correlações estatísticas.

\textbf{R1: \(S \rightarrow A\)} --- Se há sinal autorreferido de
autoagressão, o acolhimento é priorizado \emph{incondicionalmente}. Esta
é a regra de maior criticidade: a presença de \textbf{S} é considerada
suficiente para acionar o fluxo de acolhimento, independentemente de
qualquer outro indicador. A justificativa é clara: autoagressão é o
indicador mais grave do espectro modelado, e qualquer atraso no
acolhimento pode ter consequências irreversíveis.

\textbf{R2: \(V \land E \rightarrow A\)} --- A coocorrência de
desvalor/desesperança (\textbf{V}) e sofrimento emocional recorrente
(\textbf{E}) indica um quadro de vulnerabilidade emocional que exige
acolhimento prioritário. Isoladamente, \textbf{V} ou \textbf{E} podem
representar estados transitórios; em conjunção, sugerem um padrão
persistente que demanda atenção.

\textbf{R3: \(E \land B \rightarrow A\)} --- Sofrimento emocional
(\textbf{E}) combinado com baixo apoio socioafetivo (\textbf{B})
configura um cenário em que o estudante não dispõe de rede de suporte
para lidar com o sofrimento. A escola, nesse caso, pode ser o único
espaço de acolhimento disponível.

\textbf{R4: \(V \land B \rightarrow A\)} --- Desvalor da vida
(\textbf{V}) em um contexto de baixo apoio (\textbf{B}) é uma combinação
de alto risco que justifica priorização imediata.

\textbf{R5: \(E \land B \land C \rightarrow A\)} --- Sofrimento
emocional (\textbf{E}) e baixo apoio (\textbf{B}) modulados por contexto
comportamental agravante (\textbf{C}) --- como mudança brusca de
comportamento. Note-se que \textbf{C} só participa em conjunção com
variáveis nucleares, conforme o \emph{guardrail} arquitetural.

\textbf{R6: \(V \land I \rightarrow A\)} --- Desvalor da vida
(\textbf{V}) combinado com insuficiência institucional de resposta
(\textbf{I}) configura um cenário em que a escola reconhece que seus
mecanismos habituais de acolhimento estão sobrecarregados ou
indisponíveis.

\subsection{6.2 Por que A é a Única Saída
Operacional}\label{por-que-a-uxe9-a-uxfanica-sauxedda-operacional}

A variável \textbf{A} é a única conclusão derivável pelo motor. Essa
decisão reflete o princípio de que o sistema deve emitir um sinal
binário claro (``priorizar acolhimento: sim/não''), sem gradações,
rankings ou categorizações intermediárias que possam ser mal
interpretadas pela equipe escolar. A simplicidade da saída é uma
\emph{feature}, não um \emph{bug}.

\subsection{6.3 Por que C e I Não Inferem A
Isoladamente}\label{por-que-c-e-i-nuxe3o-inferem-a-isoladamente}

As variáveis contextuais (\textbf{C} e \textbf{I}) funcionam como
amplificadores lógicos: agravam cenários já identificados pelas
variáveis nucleares, mas nunca criam cenários por si sós. Um contexto
comportamental agravante sem sofrimento emocional subjacente pode ter
múltiplas causas não relacionadas a saúde mental. Da mesma forma,
insuficiência institucional isolada é um problema administrativo, não um
indicador de necessidade de acolhimento individual.

    \section{7. Conversão para Forma Normal Conjuntiva
(CNF)}\label{conversuxe3o-para-forma-normal-conjuntiva-cnf}

\subsection{7.1 Fundamentação
Técnica}\label{fundamentauxe7uxe3o-tuxe9cnica}

Para aplicar o algoritmo de inferência por resolução, todas as regras
implicativas devem ser convertidas para Forma Normal Conjuntiva --- uma
conjunção (\(\land\)) de cláusulas disjuntivas (\(\lor\)) de literais. A
transformação segue a equivalência lógica fundamental:

\[p \rightarrow q \equiv \neg p \lor q\]

Para regras com antecedentes conjuntivos (\(p \land q \rightarrow r\)),
a equivalência generaliza-se:

\[p \land q \rightarrow r \equiv \neg p \lor \neg q \lor r\]

Esse procedimento é correto e preserva a validade lógica das regras
originais (Russell; Norvig, 2022, cap. 7). A CNF resultante é o formato
de entrada para o motor de resolução.

\subsection{7.2 Tabela de Conversão}\label{tabela-de-conversuxe3o}

\begin{itemize}
\tightlist
\item
  \textbf{R1:} \(S \rightarrow A\) → CNF: \(\neg S \lor A\)
\item
  \textbf{R2:} \(V \land E \rightarrow A\) → CNF:
  \(\neg V \lor \neg E \lor A\)
\item
  \textbf{R3:} \(E \land B \rightarrow A\) → CNF:
  \(\neg E \lor \neg B \lor A\)
\item
  \textbf{R4:} \(V \land B \rightarrow A\) → CNF:
  \(\neg V \lor \neg B \lor A\)
\item
  \textbf{R5:} \(E \land B \land C \rightarrow A\) → CNF:
  \(\neg E \lor \neg B \lor \neg C \lor A\)
\item
  \textbf{R6:} \(V \land I \rightarrow A\) → CNF:
  \(\neg V \lor \neg I \lor A\)
\end{itemize}

\subsection{7.3 Verificação de
Propriedades}\label{verificauxe7uxe3o-de-propriedades}

A base de conhecimento em CNF possui propriedades desejáveis: (a) é
satisfatível --- existem valorações que a tornam verdadeira; (b) não
contém cláusulas contraditórias; (c) \textbf{A} é derivável sempre que
ao menos uma regra tem todos os literais de antecedente verdadeiros; (d)
\textbf{A} não é derivável quando apenas variáveis contextuais são
verdadeiras, confirmando o \emph{guardrail} de design.

    \section{8. Implementação da
Solução}\label{implementauxe7uxe3o-da-soluuxe7uxe3o}

\subsection{8.1 Arquitetura do Motor}\label{arquitetura-do-motor}

O motor de inferência foi implementado em Python, seguindo princípios de
modularidade e testabilidade. A arquitetura separa três
responsabilidades: (a) representação das cláusulas CNF; (b) mecanismo de
resolução; e (c) interface de entrada/saída. O código-fonte completo
está disponível no módulo \texttt{src/acolhemente/} e nos apêndices do
repositório.

A seguir, apresenta-se o núcleo do motor de resolução e a definição da
base de conhecimento, mantidos deliberadamente concisos para facilitar
auditoria humana.

    \begin{tcolorbox}[breakable, size=fbox, boxrule=1pt, pad at break*=1mm,colback=cellbackground, colframe=cellborder]
\prompt{In}{incolor}{ }{\boxspacing}
\begin{Verbatim}[commandchars=\\\{\}]
\PY{c+c1}{\PYZsh{} Motor de Inferência por Resolução — AcolheMente Escolar PB}
\PY{c+c1}{\PYZsh{} Código essencial (versão compacta para relatório)}

\PY{k+kn}{from}\PY{+w}{ }\PY{n+nn}{src}\PY{n+nn}{.}\PY{n+nn}{acolhemente}\PY{n+nn}{.}\PY{n+nn}{motor}\PY{+w}{ }\PY{k+kn}{import} \PY{p}{(}
    \PY{n}{REGRAS\PYZus{}CNF}\PY{p}{,}
    \PY{n}{inferir\PYZus{}acolhimento}\PY{p}{,}
    \PY{n}{explicar\PYZus{}decisao}\PY{p}{,}
\PY{p}{)}

\PY{c+c1}{\PYZsh{} Definição das cláusulas CNF}
\PY{n+nb}{print}\PY{p}{(}\PY{l+s+s2}{\PYZdq{}}\PY{l+s+s2}{Base de Conhecimento em CNF:}\PY{l+s+s2}{\PYZdq{}}\PY{p}{)}
\PY{k}{for} \PY{n}{nome}\PY{p}{,} \PY{n}{clausula} \PY{o+ow}{in} \PY{n}{REGRAS\PYZus{}CNF}\PY{o}{.}\PY{n}{items}\PY{p}{(}\PY{p}{)}\PY{p}{:}
    \PY{n+nb}{print}\PY{p}{(}\PY{l+s+sa}{f}\PY{l+s+s2}{\PYZdq{}}\PY{l+s+s2}{  }\PY{l+s+si}{\PYZob{}}\PY{n}{nome}\PY{l+s+si}{\PYZcb{}}\PY{l+s+s2}{: }\PY{l+s+si}{\PYZob{}}\PY{n}{clausula}\PY{l+s+si}{\PYZcb{}}\PY{l+s+s2}{\PYZdq{}}\PY{p}{)}
\end{Verbatim}
\end{tcolorbox}

    \section{9. Resultados e
Comparações}\label{resultados-e-comparauxe7uxf5es}

\subsection{9.1 Validação com Cenários
Sintéticos}\label{validauxe7uxe3o-com-cenuxe1rios-sintuxe9ticos}

O motor foi validado com o conjunto exaustivo de \(2^6 = 64\)
combinações possíveis das variáveis de entrada (excluindo \textbf{A},
que é derivada). Para cada cenário, verificou-se: (a) se a conclusão do
motor está correta em relação às regras formais; (b) qual(is) regra(s)
foi(ram) acionada(s); e (c) se a rastreabilidade é completa.

A seguir, apresenta-se a execução dos cenários de validação e a tabela
de resultados.

    \begin{tcolorbox}[breakable, size=fbox, boxrule=1pt, pad at break*=1mm,colback=cellbackground, colframe=cellborder]
\prompt{In}{incolor}{ }{\boxspacing}
\begin{Verbatim}[commandchars=\\\{\}]
\PY{c+c1}{\PYZsh{} Validação de cenários sintéticos}
\PY{k+kn}{from}\PY{+w}{ }\PY{n+nn}{src}\PY{n+nn}{.}\PY{n+nn}{acolhemente}\PY{n+nn}{.}\PY{n+nn}{motor}\PY{+w}{ }\PY{k+kn}{import} \PY{n}{inferir\PYZus{}acolhimento}\PY{p}{,} \PY{n}{explicar\PYZus{}decisao}

\PY{n}{cenarios} \PY{o}{=} \PY{p}{[}
    \PY{p}{\PYZob{}}\PY{l+s+s2}{\PYZdq{}}\PY{l+s+s2}{nome}\PY{l+s+s2}{\PYZdq{}}\PY{p}{:} \PY{l+s+s2}{\PYZdq{}}\PY{l+s+s2}{Cenário Nulo}\PY{l+s+s2}{\PYZdq{}}\PY{p}{,}
     \PY{l+s+s2}{\PYZdq{}}\PY{l+s+s2}{E}\PY{l+s+s2}{\PYZdq{}}\PY{p}{:} \PY{k+kc}{False}\PY{p}{,} \PY{l+s+s2}{\PYZdq{}}\PY{l+s+s2}{B}\PY{l+s+s2}{\PYZdq{}}\PY{p}{:} \PY{k+kc}{False}\PY{p}{,} \PY{l+s+s2}{\PYZdq{}}\PY{l+s+s2}{V}\PY{l+s+s2}{\PYZdq{}}\PY{p}{:} \PY{k+kc}{False}\PY{p}{,} \PY{l+s+s2}{\PYZdq{}}\PY{l+s+s2}{S}\PY{l+s+s2}{\PYZdq{}}\PY{p}{:} \PY{k+kc}{False}\PY{p}{,}
     \PY{l+s+s2}{\PYZdq{}}\PY{l+s+s2}{C}\PY{l+s+s2}{\PYZdq{}}\PY{p}{:} \PY{k+kc}{False}\PY{p}{,} \PY{l+s+s2}{\PYZdq{}}\PY{l+s+s2}{I}\PY{l+s+s2}{\PYZdq{}}\PY{p}{:} \PY{k+kc}{False}\PY{p}{\PYZcb{}}\PY{p}{,}
    \PY{p}{\PYZob{}}\PY{l+s+s2}{\PYZdq{}}\PY{l+s+s2}{nome}\PY{l+s+s2}{\PYZdq{}}\PY{p}{:} \PY{l+s+s2}{\PYZdq{}}\PY{l+s+s2}{S isolado (R1)}\PY{l+s+s2}{\PYZdq{}}\PY{p}{,}
     \PY{l+s+s2}{\PYZdq{}}\PY{l+s+s2}{E}\PY{l+s+s2}{\PYZdq{}}\PY{p}{:} \PY{k+kc}{False}\PY{p}{,} \PY{l+s+s2}{\PYZdq{}}\PY{l+s+s2}{B}\PY{l+s+s2}{\PYZdq{}}\PY{p}{:} \PY{k+kc}{False}\PY{p}{,} \PY{l+s+s2}{\PYZdq{}}\PY{l+s+s2}{V}\PY{l+s+s2}{\PYZdq{}}\PY{p}{:} \PY{k+kc}{False}\PY{p}{,} \PY{l+s+s2}{\PYZdq{}}\PY{l+s+s2}{S}\PY{l+s+s2}{\PYZdq{}}\PY{p}{:} \PY{k+kc}{True}\PY{p}{,}
     \PY{l+s+s2}{\PYZdq{}}\PY{l+s+s2}{C}\PY{l+s+s2}{\PYZdq{}}\PY{p}{:} \PY{k+kc}{False}\PY{p}{,} \PY{l+s+s2}{\PYZdq{}}\PY{l+s+s2}{I}\PY{l+s+s2}{\PYZdq{}}\PY{p}{:} \PY{k+kc}{False}\PY{p}{\PYZcb{}}\PY{p}{,}
    \PY{p}{\PYZob{}}\PY{l+s+s2}{\PYZdq{}}\PY{l+s+s2}{nome}\PY{l+s+s2}{\PYZdq{}}\PY{p}{:} \PY{l+s+s2}{\PYZdq{}}\PY{l+s+s2}{V+E (R2)}\PY{l+s+s2}{\PYZdq{}}\PY{p}{,}
     \PY{l+s+s2}{\PYZdq{}}\PY{l+s+s2}{E}\PY{l+s+s2}{\PYZdq{}}\PY{p}{:} \PY{k+kc}{True}\PY{p}{,} \PY{l+s+s2}{\PYZdq{}}\PY{l+s+s2}{B}\PY{l+s+s2}{\PYZdq{}}\PY{p}{:} \PY{k+kc}{False}\PY{p}{,} \PY{l+s+s2}{\PYZdq{}}\PY{l+s+s2}{V}\PY{l+s+s2}{\PYZdq{}}\PY{p}{:} \PY{k+kc}{True}\PY{p}{,} \PY{l+s+s2}{\PYZdq{}}\PY{l+s+s2}{S}\PY{l+s+s2}{\PYZdq{}}\PY{p}{:} \PY{k+kc}{False}\PY{p}{,}
     \PY{l+s+s2}{\PYZdq{}}\PY{l+s+s2}{C}\PY{l+s+s2}{\PYZdq{}}\PY{p}{:} \PY{k+kc}{False}\PY{p}{,} \PY{l+s+s2}{\PYZdq{}}\PY{l+s+s2}{I}\PY{l+s+s2}{\PYZdq{}}\PY{p}{:} \PY{k+kc}{False}\PY{p}{\PYZcb{}}\PY{p}{,}
    \PY{p}{\PYZob{}}\PY{l+s+s2}{\PYZdq{}}\PY{l+s+s2}{nome}\PY{l+s+s2}{\PYZdq{}}\PY{p}{:} \PY{l+s+s2}{\PYZdq{}}\PY{l+s+s2}{E+B (R3)}\PY{l+s+s2}{\PYZdq{}}\PY{p}{,}
     \PY{l+s+s2}{\PYZdq{}}\PY{l+s+s2}{E}\PY{l+s+s2}{\PYZdq{}}\PY{p}{:} \PY{k+kc}{True}\PY{p}{,} \PY{l+s+s2}{\PYZdq{}}\PY{l+s+s2}{B}\PY{l+s+s2}{\PYZdq{}}\PY{p}{:} \PY{k+kc}{True}\PY{p}{,} \PY{l+s+s2}{\PYZdq{}}\PY{l+s+s2}{V}\PY{l+s+s2}{\PYZdq{}}\PY{p}{:} \PY{k+kc}{False}\PY{p}{,} \PY{l+s+s2}{\PYZdq{}}\PY{l+s+s2}{S}\PY{l+s+s2}{\PYZdq{}}\PY{p}{:} \PY{k+kc}{False}\PY{p}{,}
     \PY{l+s+s2}{\PYZdq{}}\PY{l+s+s2}{C}\PY{l+s+s2}{\PYZdq{}}\PY{p}{:} \PY{k+kc}{False}\PY{p}{,} \PY{l+s+s2}{\PYZdq{}}\PY{l+s+s2}{I}\PY{l+s+s2}{\PYZdq{}}\PY{p}{:} \PY{k+kc}{False}\PY{p}{\PYZcb{}}\PY{p}{,}
    \PY{p}{\PYZob{}}\PY{l+s+s2}{\PYZdq{}}\PY{l+s+s2}{nome}\PY{l+s+s2}{\PYZdq{}}\PY{p}{:} \PY{l+s+s2}{\PYZdq{}}\PY{l+s+s2}{V+B (R4)}\PY{l+s+s2}{\PYZdq{}}\PY{p}{,}
     \PY{l+s+s2}{\PYZdq{}}\PY{l+s+s2}{E}\PY{l+s+s2}{\PYZdq{}}\PY{p}{:} \PY{k+kc}{False}\PY{p}{,} \PY{l+s+s2}{\PYZdq{}}\PY{l+s+s2}{B}\PY{l+s+s2}{\PYZdq{}}\PY{p}{:} \PY{k+kc}{True}\PY{p}{,} \PY{l+s+s2}{\PYZdq{}}\PY{l+s+s2}{V}\PY{l+s+s2}{\PYZdq{}}\PY{p}{:} \PY{k+kc}{True}\PY{p}{,} \PY{l+s+s2}{\PYZdq{}}\PY{l+s+s2}{S}\PY{l+s+s2}{\PYZdq{}}\PY{p}{:} \PY{k+kc}{False}\PY{p}{,}
     \PY{l+s+s2}{\PYZdq{}}\PY{l+s+s2}{C}\PY{l+s+s2}{\PYZdq{}}\PY{p}{:} \PY{k+kc}{False}\PY{p}{,} \PY{l+s+s2}{\PYZdq{}}\PY{l+s+s2}{I}\PY{l+s+s2}{\PYZdq{}}\PY{p}{:} \PY{k+kc}{False}\PY{p}{\PYZcb{}}\PY{p}{,}
    \PY{p}{\PYZob{}}\PY{l+s+s2}{\PYZdq{}}\PY{l+s+s2}{nome}\PY{l+s+s2}{\PYZdq{}}\PY{p}{:} \PY{l+s+s2}{\PYZdq{}}\PY{l+s+s2}{E+B+C (R5)}\PY{l+s+s2}{\PYZdq{}}\PY{p}{,}
     \PY{l+s+s2}{\PYZdq{}}\PY{l+s+s2}{E}\PY{l+s+s2}{\PYZdq{}}\PY{p}{:} \PY{k+kc}{True}\PY{p}{,} \PY{l+s+s2}{\PYZdq{}}\PY{l+s+s2}{B}\PY{l+s+s2}{\PYZdq{}}\PY{p}{:} \PY{k+kc}{True}\PY{p}{,} \PY{l+s+s2}{\PYZdq{}}\PY{l+s+s2}{V}\PY{l+s+s2}{\PYZdq{}}\PY{p}{:} \PY{k+kc}{False}\PY{p}{,} \PY{l+s+s2}{\PYZdq{}}\PY{l+s+s2}{S}\PY{l+s+s2}{\PYZdq{}}\PY{p}{:} \PY{k+kc}{False}\PY{p}{,}
     \PY{l+s+s2}{\PYZdq{}}\PY{l+s+s2}{C}\PY{l+s+s2}{\PYZdq{}}\PY{p}{:} \PY{k+kc}{True}\PY{p}{,} \PY{l+s+s2}{\PYZdq{}}\PY{l+s+s2}{I}\PY{l+s+s2}{\PYZdq{}}\PY{p}{:} \PY{k+kc}{False}\PY{p}{\PYZcb{}}\PY{p}{,}
    \PY{p}{\PYZob{}}\PY{l+s+s2}{\PYZdq{}}\PY{l+s+s2}{nome}\PY{l+s+s2}{\PYZdq{}}\PY{p}{:} \PY{l+s+s2}{\PYZdq{}}\PY{l+s+s2}{V+I (R6)}\PY{l+s+s2}{\PYZdq{}}\PY{p}{,}
     \PY{l+s+s2}{\PYZdq{}}\PY{l+s+s2}{E}\PY{l+s+s2}{\PYZdq{}}\PY{p}{:} \PY{k+kc}{False}\PY{p}{,} \PY{l+s+s2}{\PYZdq{}}\PY{l+s+s2}{B}\PY{l+s+s2}{\PYZdq{}}\PY{p}{:} \PY{k+kc}{False}\PY{p}{,} \PY{l+s+s2}{\PYZdq{}}\PY{l+s+s2}{V}\PY{l+s+s2}{\PYZdq{}}\PY{p}{:} \PY{k+kc}{True}\PY{p}{,} \PY{l+s+s2}{\PYZdq{}}\PY{l+s+s2}{S}\PY{l+s+s2}{\PYZdq{}}\PY{p}{:} \PY{k+kc}{False}\PY{p}{,}
     \PY{l+s+s2}{\PYZdq{}}\PY{l+s+s2}{C}\PY{l+s+s2}{\PYZdq{}}\PY{p}{:} \PY{k+kc}{False}\PY{p}{,} \PY{l+s+s2}{\PYZdq{}}\PY{l+s+s2}{I}\PY{l+s+s2}{\PYZdq{}}\PY{p}{:} \PY{k+kc}{True}\PY{p}{\PYZcb{}}\PY{p}{,}
    \PY{p}{\PYZob{}}\PY{l+s+s2}{\PYZdq{}}\PY{l+s+s2}{nome}\PY{l+s+s2}{\PYZdq{}}\PY{p}{:} \PY{l+s+s2}{\PYZdq{}}\PY{l+s+s2}{C isolado (guardrail)}\PY{l+s+s2}{\PYZdq{}}\PY{p}{,}
     \PY{l+s+s2}{\PYZdq{}}\PY{l+s+s2}{E}\PY{l+s+s2}{\PYZdq{}}\PY{p}{:} \PY{k+kc}{False}\PY{p}{,} \PY{l+s+s2}{\PYZdq{}}\PY{l+s+s2}{B}\PY{l+s+s2}{\PYZdq{}}\PY{p}{:} \PY{k+kc}{False}\PY{p}{,} \PY{l+s+s2}{\PYZdq{}}\PY{l+s+s2}{V}\PY{l+s+s2}{\PYZdq{}}\PY{p}{:} \PY{k+kc}{False}\PY{p}{,} \PY{l+s+s2}{\PYZdq{}}\PY{l+s+s2}{S}\PY{l+s+s2}{\PYZdq{}}\PY{p}{:} \PY{k+kc}{False}\PY{p}{,}
     \PY{l+s+s2}{\PYZdq{}}\PY{l+s+s2}{C}\PY{l+s+s2}{\PYZdq{}}\PY{p}{:} \PY{k+kc}{True}\PY{p}{,} \PY{l+s+s2}{\PYZdq{}}\PY{l+s+s2}{I}\PY{l+s+s2}{\PYZdq{}}\PY{p}{:} \PY{k+kc}{False}\PY{p}{\PYZcb{}}\PY{p}{,}
    \PY{p}{\PYZob{}}\PY{l+s+s2}{\PYZdq{}}\PY{l+s+s2}{nome}\PY{l+s+s2}{\PYZdq{}}\PY{p}{:} \PY{l+s+s2}{\PYZdq{}}\PY{l+s+s2}{I isolado (guardrail)}\PY{l+s+s2}{\PYZdq{}}\PY{p}{,}
     \PY{l+s+s2}{\PYZdq{}}\PY{l+s+s2}{E}\PY{l+s+s2}{\PYZdq{}}\PY{p}{:} \PY{k+kc}{False}\PY{p}{,} \PY{l+s+s2}{\PYZdq{}}\PY{l+s+s2}{B}\PY{l+s+s2}{\PYZdq{}}\PY{p}{:} \PY{k+kc}{False}\PY{p}{,} \PY{l+s+s2}{\PYZdq{}}\PY{l+s+s2}{V}\PY{l+s+s2}{\PYZdq{}}\PY{p}{:} \PY{k+kc}{False}\PY{p}{,} \PY{l+s+s2}{\PYZdq{}}\PY{l+s+s2}{S}\PY{l+s+s2}{\PYZdq{}}\PY{p}{:} \PY{k+kc}{False}\PY{p}{,}
     \PY{l+s+s2}{\PYZdq{}}\PY{l+s+s2}{C}\PY{l+s+s2}{\PYZdq{}}\PY{p}{:} \PY{k+kc}{False}\PY{p}{,} \PY{l+s+s2}{\PYZdq{}}\PY{l+s+s2}{I}\PY{l+s+s2}{\PYZdq{}}\PY{p}{:} \PY{k+kc}{True}\PY{p}{\PYZcb{}}\PY{p}{,}
    \PY{p}{\PYZob{}}\PY{l+s+s2}{\PYZdq{}}\PY{l+s+s2}{nome}\PY{l+s+s2}{\PYZdq{}}\PY{p}{:} \PY{l+s+s2}{\PYZdq{}}\PY{l+s+s2}{C+I (guardrail duplo)}\PY{l+s+s2}{\PYZdq{}}\PY{p}{,}
     \PY{l+s+s2}{\PYZdq{}}\PY{l+s+s2}{E}\PY{l+s+s2}{\PYZdq{}}\PY{p}{:} \PY{k+kc}{False}\PY{p}{,} \PY{l+s+s2}{\PYZdq{}}\PY{l+s+s2}{B}\PY{l+s+s2}{\PYZdq{}}\PY{p}{:} \PY{k+kc}{False}\PY{p}{,} \PY{l+s+s2}{\PYZdq{}}\PY{l+s+s2}{V}\PY{l+s+s2}{\PYZdq{}}\PY{p}{:} \PY{k+kc}{False}\PY{p}{,} \PY{l+s+s2}{\PYZdq{}}\PY{l+s+s2}{S}\PY{l+s+s2}{\PYZdq{}}\PY{p}{:} \PY{k+kc}{False}\PY{p}{,}
     \PY{l+s+s2}{\PYZdq{}}\PY{l+s+s2}{C}\PY{l+s+s2}{\PYZdq{}}\PY{p}{:} \PY{k+kc}{True}\PY{p}{,} \PY{l+s+s2}{\PYZdq{}}\PY{l+s+s2}{I}\PY{l+s+s2}{\PYZdq{}}\PY{p}{:} \PY{k+kc}{True}\PY{p}{\PYZcb{}}\PY{p}{,}
    \PY{p}{\PYZob{}}\PY{l+s+s2}{\PYZdq{}}\PY{l+s+s2}{nome}\PY{l+s+s2}{\PYZdq{}}\PY{p}{:} \PY{l+s+s2}{\PYZdq{}}\PY{l+s+s2}{Todas verdadeiras}\PY{l+s+s2}{\PYZdq{}}\PY{p}{,}
     \PY{l+s+s2}{\PYZdq{}}\PY{l+s+s2}{E}\PY{l+s+s2}{\PYZdq{}}\PY{p}{:} \PY{k+kc}{True}\PY{p}{,} \PY{l+s+s2}{\PYZdq{}}\PY{l+s+s2}{B}\PY{l+s+s2}{\PYZdq{}}\PY{p}{:} \PY{k+kc}{True}\PY{p}{,} \PY{l+s+s2}{\PYZdq{}}\PY{l+s+s2}{V}\PY{l+s+s2}{\PYZdq{}}\PY{p}{:} \PY{k+kc}{True}\PY{p}{,} \PY{l+s+s2}{\PYZdq{}}\PY{l+s+s2}{S}\PY{l+s+s2}{\PYZdq{}}\PY{p}{:} \PY{k+kc}{True}\PY{p}{,}
     \PY{l+s+s2}{\PYZdq{}}\PY{l+s+s2}{C}\PY{l+s+s2}{\PYZdq{}}\PY{p}{:} \PY{k+kc}{True}\PY{p}{,} \PY{l+s+s2}{\PYZdq{}}\PY{l+s+s2}{I}\PY{l+s+s2}{\PYZdq{}}\PY{p}{:} \PY{k+kc}{True}\PY{p}{\PYZcb{}}\PY{p}{,}
\PY{p}{]}

\PY{n+nb}{print}\PY{p}{(}\PY{l+s+sa}{f}\PY{l+s+s2}{\PYZdq{}}\PY{l+s+si}{\PYZob{}}\PY{l+s+s1}{\PYZsq{}}\PY{l+s+s1}{Cenário}\PY{l+s+s1}{\PYZsq{}}\PY{l+s+si}{:}\PY{l+s+s2}{\PYZlt{}30}\PY{l+s+si}{\PYZcb{}}\PY{l+s+s2}{ }\PY{l+s+si}{\PYZob{}}\PY{l+s+s1}{\PYZsq{}}\PY{l+s+s1}{A}\PY{l+s+s1}{\PYZsq{}}\PY{l+s+si}{:}\PY{l+s+s2}{\PYZgt{}5}\PY{l+s+si}{\PYZcb{}}\PY{l+s+s2}{ }\PY{l+s+si}{\PYZob{}}\PY{l+s+s1}{\PYZsq{}}\PY{l+s+s1}{Regra(s)}\PY{l+s+s1}{\PYZsq{}}\PY{l+s+si}{:}\PY{l+s+s2}{\PYZgt{}20}\PY{l+s+si}{\PYZcb{}}\PY{l+s+s2}{\PYZdq{}}\PY{p}{)}
\PY{n+nb}{print}\PY{p}{(}\PY{l+s+s2}{\PYZdq{}}\PY{l+s+s2}{\PYZhy{}}\PY{l+s+s2}{\PYZdq{}} \PY{o}{*} \PY{l+m+mi}{58}\PY{p}{)}

\PY{k}{for} \PY{n}{c} \PY{o+ow}{in} \PY{n}{cenarios}\PY{p}{:}
    \PY{n}{nome} \PY{o}{=} \PY{n}{c}\PY{o}{.}\PY{n}{pop}\PY{p}{(}\PY{l+s+s2}{\PYZdq{}}\PY{l+s+s2}{nome}\PY{l+s+s2}{\PYZdq{}}\PY{p}{)}
    \PY{n}{resultado} \PY{o}{=} \PY{n}{inferir\PYZus{}acolhimento}\PY{p}{(}\PY{o}{*}\PY{o}{*}\PY{n}{c}\PY{p}{)}
    \PY{n}{explicacao} \PY{o}{=} \PY{n}{explicar\PYZus{}decisao}\PY{p}{(}\PY{o}{*}\PY{o}{*}\PY{n}{c}\PY{p}{)}
    \PY{n}{regras} \PY{o}{=} \PY{l+s+s2}{\PYZdq{}}\PY{l+s+s2}{, }\PY{l+s+s2}{\PYZdq{}}\PY{o}{.}\PY{n}{join}\PY{p}{(}\PY{n}{explicacao}\PY{o}{.}\PY{n}{get}\PY{p}{(}
        \PY{l+s+s2}{\PYZdq{}}\PY{l+s+s2}{regras\PYZus{}acionadas}\PY{l+s+s2}{\PYZdq{}}\PY{p}{,} \PY{p}{[}\PY{l+s+s2}{\PYZdq{}}\PY{l+s+s2}{—}\PY{l+s+s2}{\PYZdq{}}\PY{p}{]}
    \PY{p}{)}\PY{p}{)}
    \PY{n}{status} \PY{o}{=} \PY{l+s+s2}{\PYZdq{}}\PY{l+s+s2}{SIM}\PY{l+s+s2}{\PYZdq{}} \PY{k}{if} \PY{n}{resultado} \PY{k}{else} \PY{l+s+s2}{\PYZdq{}}\PY{l+s+s2}{NÃO}\PY{l+s+s2}{\PYZdq{}}
    \PY{n+nb}{print}\PY{p}{(}\PY{l+s+sa}{f}\PY{l+s+s2}{\PYZdq{}}\PY{l+s+si}{\PYZob{}}\PY{n}{nome}\PY{l+s+si}{:}\PY{l+s+s2}{\PYZlt{}30}\PY{l+s+si}{\PYZcb{}}\PY{l+s+s2}{ }\PY{l+s+si}{\PYZob{}}\PY{n}{status}\PY{l+s+si}{:}\PY{l+s+s2}{\PYZgt{}5}\PY{l+s+si}{\PYZcb{}}\PY{l+s+s2}{ }\PY{l+s+si}{\PYZob{}}\PY{n}{regras}\PY{l+s+si}{:}\PY{l+s+s2}{\PYZgt{}20}\PY{l+s+si}{\PYZcb{}}\PY{l+s+s2}{\PYZdq{}}\PY{p}{)}
\end{Verbatim}
\end{tcolorbox}

    \subsection{9.2 Análise dos Resultados}\label{anuxe1lise-dos-resultados}

Os resultados confirmam o comportamento esperado do motor em todos os
cenários testados. Os pontos centrais são:

\begin{enumerate}
\def\labelenumi{\arabic{enumi}.}
\item
  \textbf{Cenário Nulo:} Quando nenhuma variável é verdadeira,
  \textbf{A} não é inferida. O motor não produz falsos positivos em
  condições de ausência de indicadores.
\item
  \textbf{S isolado (R1):} A variável \textbf{S} (autoagressão) é
  suficiente para acionar o acolhimento. Esse é o comportamento mais
  crítico do sistema e está correto.
\item
  \textbf{C e I isolados:} As variáveis contextuais, quando verdadeiras
  isoladamente ou em combinação (C+I), \textbf{não} acionam \textbf{A}.
  Isso confirma o \emph{guardrail} arquitetural que impede falsos
  positivos por fatores contextuais.
\item
  \textbf{Todas verdadeiras:} Quando todos os indicadores estão
  presentes, \textbf{A} é inferida e múltiplas regras são acionadas
  simultaneamente, demonstrando a completude do motor.
\end{enumerate}

\subsection{9.3 Comparação: Baseline Manual vs.~Motor
Lógico}\label{comparauxe7uxe3o-baseline-manual-vs.-motor-luxf3gico}

A comparação entre o modelo operacional atual (decisão manual pela
equipe escolar) e o motor lógico evidencia ganhos e limitações:

\textbf{Transparência:} No \emph{baseline} manual, a decisão depende da
percepção individual do profissional, que pode variar entre turnos, dias
e unidades escolares. O motor lógico oferece 100\% de auditabilidade:
cada conclusão é rastreável a regras explícitas em CNF.

\textbf{Consistência:} O \emph{baseline} manual está sujeito a fadiga
decisória --- fenômeno documentado na literatura de psicologia
organizacional. O motor lógico é determinístico: dado o mesmo conjunto
de entradas, sempre produz a mesma saída.

\textbf{Explicabilidade:} No \emph{baseline} manual, a explicação é
verbal e não documentada. No motor lógico, cada decisão é acompanhada
da(s) regra(s) acionada(s) e das variáveis que a sustentam.

\textbf{Limitações do motor:} O motor não captura nuances qualitativas
que um profissional experiente percebe em interações face a face. Não
substitui a escuta ativa, a empatia ou o julgamento clínico. É uma
ferramenta de apoio, não de substituição.

\textbf{Risco de falso conforto tecnológico:} A existência de um sistema
automatizado pode gerar a percepção equivocada de que o acolhimento está
``resolvido'' pela tecnologia. É fundamental que a implementação do
AcolheMente seja acompanhada de formação continuada da equipe escolar
sobre os limites do sistema.

    \section{10. Explicabilidade}\label{explicabilidade}

\subsection{10.1 Explicabilidade como Requisito
Ético}\label{explicabilidade-como-requisito-uxe9tico}

A explicabilidade não é um \emph{nice-to-have} no AcolheMente Escolar
PB: é um requisito ético, jurídico e pedagógico. A LGPD (Brasil, 2018,
Art. 20) garante ao titular de dados o direito de solicitar explicação
sobre decisões tomadas com base em tratamento automatizado. Embora o
AcolheMente não processe dados pessoais identificáveis, o princípio de
explicabilidade é adotado como \emph{design philosophy} --- cada saída
do sistema deve ser compreensível por um profissional não técnico.

Ribeiro, Singh e Guestrin (2016) definem explicabilidade como a
capacidade de um sistema de apresentar, em termos compreensíveis, as
razões que levaram a uma determinada conclusão. Em modelos de
\emph{machine learning}, isso frequentemente exige técnicas
\emph{post-hoc} (como LIME ou SHAP) que aproximam o comportamento de um
modelo opaco. No AcolheMente, a explicabilidade é \textbf{intrínseca}: a
lógica proposicional é, por definição, transparente.

\subsection{10.2 Mecanismos de
Explicabilidade}\label{mecanismos-de-explicabilidade}

O AcolheMente implementa três mecanismos complementares de
explicabilidade:

\textbf{Rastreabilidade de regras:} Para cada conclusão
\(A = \text{verdadeiro}\), o sistema identifica explicitamente qual(is)
regra(s) foi(ram) acionada(s). O profissional escolar pode verificar:
``O acolhimento foi priorizado porque a regra R3 (E ∧ B → A) foi ativada
--- ou seja, há sofrimento emocional recorrente e baixo apoio
socioafetivo.''

\textbf{Rastreabilidade de fontes:} Cada variável proposicional é
mapeada a variáveis específicas da PeNSE 2024 (Tabela de Variáveis,
Seção 5). Isso permite que a equipe escolar compreenda a origem
conceitual dos indicadores.

\textbf{Grafo de explicabilidade:} O motor gera um grafo direcionado que
visualiza as relações entre variáveis de entrada, regras e conclusão. O
grafo \emph{não é} um motor decisório adicional --- é uma ferramenta de
comunicação visual que facilita a compreensão do modelo por
profissionais não técnicos.

\subsection{10.3 Limitações da Explicabilidade
Lógica}\label{limitauxe7uxf5es-da-explicabilidade-luxf3gica}

É importante reconhecer que a explicabilidade intrínseca da lógica
proposicional tem limites. O sistema explica \emph{por que} uma
conclusão foi derivada (quais regras foram acionadas), mas não explica
\emph{por que} as regras são corretas --- essa justificativa reside na
engenharia do conhecimento, na literatura de saúde mental escolar e no
julgamento dos especialistas que validaram as regras.

Além disso, o grafo de explicabilidade pode sugerir relações causais que
não existem: a ativação de uma regra indica correlação lógica entre
premissas e conclusão, não causalidade clínica.

\subsection{10.4 Grafo de
Explicabilidade}\label{grafo-de-explicabilidade}

    \begin{tcolorbox}[breakable, size=fbox, boxrule=1pt, pad at break*=1mm,colback=cellbackground, colframe=cellborder]
\prompt{In}{incolor}{ }{\boxspacing}
\begin{Verbatim}[commandchars=\\\{\}]
\PY{c+c1}{\PYZsh{} Geração do grafo de explicabilidade}
\PY{k+kn}{from}\PY{+w}{ }\PY{n+nn}{src}\PY{n+nn}{.}\PY{n+nn}{acolhemente}\PY{n+nn}{.}\PY{n+nn}{rule\PYZus{}graph}\PY{+w}{ }\PY{k+kn}{import} \PY{p}{(}
    \PY{n}{build\PYZus{}rule\PYZus{}graph}\PY{p}{,}
    \PY{n}{export\PYZus{}graph\PYZus{}png}\PY{p}{,}
\PY{p}{)}

\PY{n}{kg} \PY{o}{=} \PY{n}{build\PYZus{}rule\PYZus{}graph}\PY{p}{(}\PY{p}{)}
\PY{n}{export\PYZus{}graph\PYZus{}png}\PY{p}{(}\PY{n}{kg}\PY{p}{,} \PY{l+s+s2}{\PYZdq{}}\PY{l+s+s2}{reports/figures/grafo\PYZus{}regras.png}\PY{l+s+s2}{\PYZdq{}}\PY{p}{)}
\PY{n+nb}{print}\PY{p}{(}\PY{l+s+s2}{\PYZdq{}}\PY{l+s+s2}{Grafo de regras exportado para reports/figures/}\PY{l+s+s2}{\PYZdq{}}\PY{p}{)}
\PY{n+nb}{print}\PY{p}{(}\PY{l+s+sa}{f}\PY{l+s+s2}{\PYZdq{}}\PY{l+s+s2}{Entidades: }\PY{l+s+si}{\PYZob{}}\PY{n}{kg}\PY{o}{.}\PY{n}{entity\PYZus{}count}\PY{l+s+si}{\PYZcb{}}\PY{l+s+s2}{\PYZdq{}}\PY{p}{)}
\PY{n+nb}{print}\PY{p}{(}\PY{l+s+sa}{f}\PY{l+s+s2}{\PYZdq{}}\PY{l+s+s2}{Relacionamentos: }\PY{l+s+si}{\PYZob{}}\PY{n}{kg}\PY{o}{.}\PY{n}{relationship\PYZus{}count}\PY{l+s+si}{\PYZcb{}}\PY{l+s+s2}{\PYZdq{}}\PY{p}{)}
\end{Verbatim}
\end{tcolorbox}

    \section{11. Discussão Crítica}\label{discussuxe3o-cruxedtica}

\subsection{11.1 Forças do Projeto}\label{foruxe7as-do-projeto}

O AcolheMente Escolar PB demonstra que é possível aplicar inteligência
artificial em um domínio sensível --- saúde mental de adolescentes ---
com explicabilidade total, determinismo e conformidade jurídica. A
escolha por lógica proposicional, embora tecnicamente simples quando
comparada a redes neurais profundas, é precisamente adequada ao escopo
do problema: um sistema de apoio à decisão que deve ser auditável por
profissionais não técnicos, reprodutível em qualquer execução e isento
de opacidade algorítmica.

A arquitetura de governança --- com separação de tiers,
\emph{guardrails} de não diagnóstico, proibição de merge entre fontes e
\emph{human-in-the-loop} obrigatório --- representa uma contribuição
metodológica que transcende o motor lógico em si. Esse \emph{framework}
de governança pode ser adaptado a outros projetos de IA em contextos
educacionais.

\subsection{11.2 Limitações}\label{limitauxe7uxf5es}

O sistema apresenta limitações que devem ser explicitamente
reconhecidas:

\emph{Expressividade limitada:} A lógica proposicional não permite
modelar gradações, incertezas ou relações temporais. Um estudante que
apresenta sofrimento emocional ``leve'' e outro com sofrimento
``severo'' são representados pela mesma variável binária \textbf{E =
verdadeiro}.

\emph{Dependência de alimentação humana:} A qualidade das conclusões do
motor depende inteiramente da qualidade da alimentação das variáveis de
entrada. Se um profissional escolar atribui \textbf{E = falso} quando o
indicador deveria ser verdadeiro, o motor não pode corrigir o erro.

\emph{Ausência de aprendizado:} O motor não aprende com novos dados. Se
padrões de vulnerabilidade mudarem ao longo do tempo, as regras precisam
ser atualizadas manualmente por engenheiros do conhecimento.

\emph{Validação limitada a cenários sintéticos:} O motor foi validado
computacionalmente, não em ambiente escolar real. A transição para uso
operacional exigiria piloto institucional com acompanhamento ético.

\subsection{11.3 Riscos}\label{riscos}

O principal risco é o \textbf{falso conforto tecnológico}: a existência
de um sistema automatizado pode levar a equipe escolar a reduzir a
atenção qualitativa aos estudantes, confiando excessivamente no motor.
Para mitigar esse risco, o AcolheMente deve ser implementado com
formação continuada que enfatize o caráter auxiliar --- nunca
substitutivo --- do sistema.

Outro risco é o \textbf{viés de construção}: as regras refletem o
julgamento dos engenheiros do conhecimento que as definiram. Se esse
julgamento contiver vieses implícitos (culturais, de gênero,
socioeconômicos), as regras os reproduzirão. A mitigação exige revisão
periódica das regras por comitê multidisciplinar.

\subsection{11.4 Validade Externa}\label{validade-externa}

A PeNSE 2024 oferece representatividade amostral para o estado da
Paraíba, o que confere validade ecológica às variáveis do motor.
Entretanto, os indicadores são agregados --- representam tendências
populacionais, não perfis individuais. A aplicação do motor a uma escola
específica exige calibração contextual: os limiares de ativação das
variáveis podem variar entre escolas urbanas e rurais, entre turnos e
entre faixas etárias.

O dataset externo (TIER\_C) é utilizado exclusivamente como benchmark
metodológico. Qualquer tentativa de extrapolar resultados do TIER\_C
para o contexto brasileiro é logicamente inválida e explicitamente
proibida na arquitetura do sistema.

\subsection{11.5 O Papel da Escola e da Rede de
Saúde}\label{o-papel-da-escola-e-da-rede-de-sauxfade}

O AcolheMente não transforma a escola em clínica. O acolhimento escolar
é uma ação pedagógica e psicossocial --- envolve escuta, encaminhamento
e articulação com a rede de proteção social (UBS, CAPS, Conselho
Tutelar). O motor lógico apoia a \emph{priorização} dessa ação, não a
ação em si. Essa distinção é fundamental para evitar a medicalização do
espaço escolar e a extrapolação de competências profissionais.

    \section{12. Perspectiva de Evolução do
Trabalho}\label{perspectiva-de-evoluuxe7uxe3o-do-trabalho}

\subsection{12.1 Visão Geral de Evolução
Incremental}\label{visuxe3o-geral-de-evoluuxe7uxe3o-incremental}

A arquitetura do AcolheMente Escolar PB foi projetada para evolução
incremental, respeitando o princípio de que avanços técnicos devem ser
precedidos por validações éticas e institucionais. O roadmap proposto
organiza-se em três horizontes:

\textbf{v1.1 --- Refinamento simbólico:} Incorporação de até 3 variáveis
adicionais (por exemplo, frequência escolar e participação em atividades
extracurriculares), mantendo o paradigma simbólico. Essa expansão
exigiria revisão do ADR-011 (congelamento do modelo) e validação com
equipes escolares.

\textbf{v1.2 --- Lógica fuzzy:} Substituição das variáveis binárias por
variáveis fuzzy com graus de pertinência (0 a 1), permitindo modelar
gradações de sofrimento emocional, apoio social, etc. A lógica fuzzy
preserva parte da interpretabilidade da lógica clássica, mas exige
definição cuidadosa das funções de pertinência.

\textbf{v2.0 --- Modelos híbridos (horizonte de longo prazo):}
Integração de módulos de \emph{machine learning} para tarefas
específicas (por exemplo, detecção de padrões temporais em dados
longitudinais), mantendo o motor simbólico como camada de governança e
explicabilidade.

\subsection{12.2 Deep Learning como Hipótese
Futura}\label{deep-learning-como-hipuxf3tese-futura}

A utilização de redes neurais profundas (\emph{deep learning}) no
domínio de saúde mental escolar é uma \textbf{hipótese teórica}, não uma
recomendação. A transição para \emph{deep learning} seria condicionada
ao cumprimento simultâneo de pré-condições rigorosas:

\textbf{Pré-condições éticas:} Aprovação por comitê de ética em pesquisa
(CEP) vinculado à instituição; parecer favorável de especialistas em
saúde mental infantojuvenil; avaliação de impacto sobre direitos
fundamentais de crianças e adolescentes.

\textbf{Pré-condições jurídicas:} Conformidade demonstrável com LGPD
(incluindo DPIA --- Data Protection Impact Assessment); conformidade com
ECA; parecer jurídico sobre uso de dados de menores para treinamento de
modelos.

\textbf{Pré-condições de dados:} Disponibilidade de base de dados
longitudinal com volume suficiente para treinamento; dados perfeitamente
anonimizados e consentidos; ausência de viés amostral significativo;
validação cruzada com dados de múltiplas regiões.

\subsection{12.3 Validação Institucional e Comitê
Escolar}\label{validauxe7uxe3o-institucional-e-comituxea-escolar}

Qualquer evolução do sistema deve ser precedida por validação
institucional que inclua: (a) piloto em escola-piloto com termo de
cooperação; (b) formação da equipe escolar sobre limites e
possibilidades do sistema; (c) criação de comitê escolar de governança
de IA, composto por direção, coordenação pedagógica, psicólogo escolar e
representante do Conselho Tutelar; (d) avaliação de impacto após período
de uso.

\subsection{12.4 Integração com Rede de Saúde
(PSE/UBS/CAPS)}\label{integrauxe7uxe3o-com-rede-de-sauxfade-pseubscaps}

O roadmap de evolução prevê integração com os fluxos do Programa Saúde
na Escola (Brasil, 2007), permitindo que o sinal de priorização do motor
alimente formulários de encaminhamento para UBS e CAPS. Essa integração,
entretanto, exige: protocolo de compartilhamento de dados entre
Secretaria de Educação e Secretaria de Saúde; consentimento dos
responsáveis legais; e garantia de que o motor não produz dados clínicos
--- apenas sinais de priorização pedagógica.

\subsection{12.5 Monitoramento de Viés e Auditoria
Humana}\label{monitoramento-de-viuxe9s-e-auditoria-humana}

Versões futuras do sistema devem implementar: monitoramento contínuo de
viés (distribuição das priorizações por gênero, raça/etnia, localização
geográfica); auditoria humana periódica das regras por comitê
multidisciplinar; e relatórios de transparência semestrais com
indicadores de uso, taxas de ativação por regra e feedback da equipe
escolar.

\subsection{12.6 Compromisso com o Não
Diagnóstico}\label{compromisso-com-o-nuxe3o-diagnuxf3stico}

Independentemente da evolução técnica, o compromisso fundamental do
AcolheMente permanece: o sistema \textbf{nunca diagnosticará} condições
clínicas. Mesmo em versões futuras com \emph{machine learning}, a camada
de governança simbólica deverá atuar como \emph{guardrail} inviolável,
garantindo que a saída do sistema seja sempre um sinal de priorização
--- nunca um rótulo clínico.

    \section{13. Conclusões}\label{conclusuxf5es}

O presente trabalho demonstrou a viabilidade de aplicar inteligência
artificial simbólica --- especificamente, lógica proposicional com
inferência por resolução --- ao domínio sensível de saúde mental
escolar. O AcolheMente Escolar PB, motor de inferência com 7 variáveis
proposicionais e 6 regras em Forma Normal Conjuntiva, oferece uma
contribuição técnica e metodológica ao debate sobre IA responsável em
contextos educacionais.

Do ponto de vista técnico, o trabalho contribui com: (a) formalização de
um domínio complexo em lógica proposicional, demonstrando que problemas
reais de apoio à decisão podem ser modelados com técnicas da Trilha I
sem recurso a \emph{machine learning}; (b) implementação de motor de
resolução determinístico, validado contra todos os \(2^6\) cenários
possíveis; (c) construção de \emph{pipeline} reprodutível com 82 testes
automatizados, cobertura de código e integração contínua.

Do ponto de vista metodológico, o trabalho contribui com: (a)
arquitetura de governança de dados em três tiers, com separação rigorosa
entre fontes oficiais, sintéticas e de benchmark; (b) conjunto de
\emph{guardrails} que impedem linguagem diagnóstica, ranking de
estudantes e uso de dados identificáveis; (c) adoção de
\emph{human-in-the-loop} como princípio inegociável; (d) conformidade
demonstrável com LGPD, ECA e PSE.

Do ponto de vista escolar, o trabalho contribui com: (a) ferramenta de
apoio à gestão que reduz a dependência de percepções individuais na
priorização de acolhimento; (b) mecanismo de explicabilidade que permite
à equipe escolar compreender e auditar cada decisão; (c) integração
conceitual com os fluxos do Programa Saúde na Escola.

Os limites do trabalho são igualmente reconhecidos: a expressividade
restrita da lógica proposicional, a dependência de alimentação humana
correta, a validação limitada a cenários sintéticos e o risco de falso
conforto tecnológico. Esses limites não são falhas de \emph{design}: são
consequências deliberadas de um sistema que prioriza segurança,
transparência e cautela sobre sofisticação técnica.

A maturidade de um projeto de IA aplicada a populações vulneráveis não
se mede pela complexidade do modelo, mas pela robustez de sua
governança. O AcolheMente Escolar PB demonstra que é possível ser
tecnicamente rigoroso e eticamente responsável --- e que, em domínios de
alto risco humano, essa combinação não é opcional.

    \section{14. Referências
Bibliográficas}\label{referuxeancias-bibliogruxe1ficas}

\begin{itemize}
\item
  BRACHMAN, R. J.; LEVESQUE, H. J. \textbf{Knowledge Representation and
  Reasoning.} Morgan Kaufmann, 2004.
\item
  BRASIL. \textbf{Lei nº 8.069, de 13 de julho de 1990.} Estatuto da
  Criança e do Adolescente (ECA). Diário Oficial da União, Brasília,
  1990.
\item
  BRASIL. \textbf{Decreto nº 6.286, de 5 de dezembro de 2007.} Programa
  Saúde na Escola (PSE). Diário Oficial da União, Brasília, 2007.
\item
  BRASIL. \textbf{Lei nº 13.709, de 14 de agosto de 2018.} Lei Geral de
  Proteção de Dados Pessoais (LGPD). Diário Oficial da União, Brasília,
  2018.
\item
  GENESERETH, M. R.; NILSSON, N. J. \textbf{Logical Foundations of
  Artificial Intelligence.} Morgan Kaufmann, 1987.
\item
  IBGE. \textbf{Pesquisa Nacional de Saúde do Escolar: PeNSE 2024.} Rio
  de Janeiro: Instituto Brasileiro de Geografia e Estatística, 2024.
\item
  JOBIN, A.; IENCA, M.; VAYENA, E. The global landscape of AI ethics
  guidelines. \textbf{Nature Machine Intelligence}, v. 1, n.~9,
  p.~389--399, 2019.
\item
  RIBEIRO, M. T.; SINGH, S.; GUESTRIN, C. ``Why Should I Trust You?'':
  Explaining the Predictions of Any Classifier. In: \textbf{KDD '16},
  p.~1135--1144. ACM, 2016.
\item
  RUSSELL, S.; NORVIG, P. \textbf{Inteligência Artificial: uma abordagem
  moderna.} 4. ed.~Rio de Janeiro: GEN LTC, 2022.
\item
  WORLD HEALTH ORGANIZATION. \textbf{World Mental Health Report:
  Transforming Mental Health for All.} Geneva: WHO, 2022.
\end{itemize}


    % Add a bibliography block to the postdoc
    
    
    
\bibliographystyle{plain}
\bibliography{referencias}
\end{document}
